\documentclass[10pt]{article}
\usepackage[margin=1in]{geometry}
\usepackage[T1]{fontenc}
\usepackage{lmodern}
\usepackage{amsmath,amssymb}
\usepackage{xcolor}
\usepackage{listings}
\usepackage{booktabs}
\usepackage{tabularx}
\usepackage{longtable}
\usepackage{pdflscape}
\usepackage{siunitx}
\usepackage{graphicx}
\usepackage[numbers,sort&compress]{natbib}
\usepackage{hyperref}
\hypersetup{colorlinks=true,linkcolor=black,urlcolor=blue,citecolor=black}

\newcommand{\ESALCI}{ESALC-I}
\newcommand{\Bmat}{\mathbf{B}}
\newcommand{\q}{\mathbf{q}}
\renewcommand{\thesection}{S\arabic{section}}
\renewcommand{\thetable}{S\arabic{table}}
\renewcommand{\thefigure}{S\arabic{figure}}
\AtBeginDocument{\renewcommand{\thelstlisting}{S\arabic{lstlisting}}}
\makeatletter
\renewcommand*\l@section{\@dottedtocline{1}{0em}{3.2em}}
\renewcommand*\l@subsection{\@dottedtocline{2}{3.2em}{4.2em}}
\renewcommand*\l@subsubsection{\@dottedtocline{3}{7.4em}{5.2em}}
\makeatother

\lstset{
  basicstyle=\ttfamily\scriptsize,
  breaklines=true,
  breakatwhitespace=false,
  columns=fullflexible,
  keepspaces=true,
  frame=single,
  rulecolor=\color{black!35},
  xleftmargin=0.5em,
  xrightmargin=0.5em
}

\title{Supporting Information for\\
Local, Non-Redundant Molecular Coordinates from Disjoint Primitive Families: Symmetry-Adapted Construction and Applications}
\author{Vincenzo Barone}
\date{}

\begin{document}
\maketitle
\tableofcontents

{This Supporting Information (SI) documents extended symmetry-adapted linear combinations
(ESALCs) in internal (ESALC-I) and Cartesian (ESALC-C) form. The extension
combines chemically assigned primitive families, local symmetry-adapted
combinations, exact-rank selection, and molecular-symmetry orientation. This
SI documents their construction, structural validation, and use in external
optimization programs.  Continuous chemical perception supplies the molecular state used for
coordinate construction.\cite{Barone2026Perception}}
{Gaussian\cite{GaussianGICManual} performs the geometry
optimizations, and the \href{https://pyscf.org/}{Python-based Simulations
of Chemistry Framework (PySCF)} supplies electronic energies and gradients
for the common-potential comparisons.}

\section{Molecular State and Coordinate Construction}

\subsection{Chemical-domain assignment}

{ESALC-I assigns a coordinate language to each chemical domain,
closes the physical rank, and applies molecular symmetry within each family.
The primitive pool supports this construction; the final set is non-redundant.
Table~\ref{tab:si-domain-policy} summarizes the
domain choices.}

\begin{table}[htbp]
\centering
\small
\caption{{Coordinate descriptions assigned to chemical domains.}}
\label{tab:si-domain-policy}
\begin{tabularx}{\linewidth}{@{}>{\raggedright\arraybackslash}p{0.21\linewidth}>{\raggedright\arraybackslash}p{0.36\linewidth}>{\raggedright\arraybackslash}X@{}}
\toprule
Chemical domain & Constructed description & Reason for the choice \\
\midrule
Covalent region & Local natural-coordinate-like blocks with exact-rank closure
& Preserves chemical locality while avoiding null directions. \\
Whole fragments & Intrafragment coordinates plus relative center translations
and quaternion-derived exponential-map orientations & Represents rigid fragment motion without
inventing interfragment bonds. \\
Rings and interaction centers & Ring deformation and puckering; atom--center
coordinates for fragment, ring, and \(\eta^n\) centers & Retains collective
motions and non-atomic interaction sites explicitly. \\
High coordination & Delocalized \(1/R\) combinations reduced to exact rank and
symmetry oriented & Avoids an arbitrary valence graph for cluster-like
environments. \\
\bottomrule
\end{tabularx}
\end{table}

{A basin is a region in which the assigned topology, primitive
identities, and selected molecular point group remain fixed. The group is the
highest symmetry intended to be conserved throughout the basin; it can be
specified explicitly with \texttt{symmetry\_group=...}. It may be a subgroup
of the initial geometry's full point group. Chemical perception
identifies domains; primitive functions
are assigned to these domains, combined into local symmetry-adapted linear
combinations (SALCs), reduced by singular-value analysis of the Wilson B matrix,
and oriented under molecular symmetry.  Primitive identities and the number of
ESALC-I coordinates remain fixed within that basin.  Single-atom out-of-plane
(SOOP) functions are linear
Cartesian combinations whose coefficients follow geometry without changing
the coordinate identities. At saturated tetracoordinate ring atoms,
only ring-owned endocyclic SOOP directions enter the puckering block; valence
angles span the other local angular motions, so no separate umbrella
coordinate is assigned. Poor conditioning or ambiguous primitives prompt
reassessment of the molecular state.  When the primitive assignment becomes unreliable,
a complete ESALC-C basis preserves the molecular symmetry and vibrational rank.
Both coordinate forms extend the SALC construction to complete,
non-redundant bases of local source directions.
}

\subsection{Fixed molecular state and Gaussian input}

The perception procedure supplies the molecular topology, ring and
multicenter domains, fragments, and symmetry operations. It uses continuous measures of bonding,
coordination, and local equivalence to identify these domains. Primitive
functions and their family assignments remain fixed in a well-defined
basin; a change in topology or retained symmetry, or an ambiguous primitive
assignment, calls for a new construction.

The coordinate builder writes the complete ESALC-I set in a
\texttt{ReadAllGIC} input for Gaussian 16 Rev. C.02. The input
includes independent coordinates and frozen symmetry partners. The
construction and Gaussian optimization therefore use the same coordinate
definitions.

For disconnected fragments, a reference fragment defines relative
translations and orientations. Open contacts use local interaction
distances and orientations; closed contact cycles use relative-fragment
motion. Shared-proton and group-13 hydride bridges are treated as
multicenter chemical domains. Each description retains the full
non-redundant rank.

\subsection{{Coordinate-builder input formats}}

{The builder accepts a Cartesian coordinate file in XYZ format (atom count, comment line, and
element symbols with Cartesian coordinates in angstrom) followed, when
needed, by a \texttt{\#SMITH} section. The input format identifier is
\nolinkurl{matrix.smith.input.v1}. Construction options in that section can
specify the retained symmetry group and fragment treatment.
Table~\ref{tab:si-builder-input} distinguishes the four input modes by the
molecular information supplied with the geometry.}

{
\begin{table}[htbp]
\centering
\small
\caption{{Molecular information accepted by the coordinate builder.}}
\label{tab:si-builder-input}
\begin{tabularx}{\linewidth}{@{}>{\raggedright\arraybackslash}p{0.26\linewidth}X@{}}
\toprule
Input mode & Supplied information and construction step \\
\midrule
Complete molecular state & Topology, primitive coordinates, molecular-symmetry
operations and atom permutations are supplied and retained. Ring, fragment,
interaction-center, and continuous-perception information may also be supplied. \\
Supplied topology & A \texttt{TOPOLOGY} section specifies covalent connectivity;
ordinary redundant primitives are then generated from it. \\
Supplied primitives & A \texttt{PRIMITIVES} section specifies the ordered
redundant coordinate functions; their atom indices, values, and analytic Wilson
B matrix rows are checked against the geometry. \\
Cartesian geometry & The included perception procedure supplies ordinary
topology, molecular symmetry with atom permutations, and primitive coordinates. \\
\bottomrule
\end{tabularx}
\end{table}
}

{When topology and primitives are both supplied, their bond
definitions must agree. A complete molecular state is useful when continuous
chemical descriptors or special interaction centers have already been
established by chemical perception. The selected input mode, accepted coordinate options,
primitive definitions, and symmetry operations accompany the constructed
coordinate set, so its Wilson B matrix can be evaluated at later geometries
within the same basin.}

{A minimal water input illustrates the Cartesian format and
the construction section:}
\begin{lstlisting}[basicstyle=\ttfamily\scriptsize]
3
Water
O  0.00000000  0.00000000  0.00000000
H  0.75860200  0.00000000  0.50428400
H -0.75860200  0.00000000  0.50428400

#SMITH
schema = matrix.smith.input.v1
symmetrize = true
\end{lstlisting}


{A complete norbornane example gives the starting geometry,
construction directives, and the generated Gaussian 16 Rev. C.02
\texttt{ReadAllGIC} input. The latter contains all 48 GIC expressions and ten
component dihedrals used in this export.}

\begin{lstlisting}[basicstyle=\ttfamily\scriptsize,caption={Complete norbornane construction input.}]
19
Norbornane starting geometry for ESALC-I ReadAllGIC export
C       1.25441300      0.78233900     -0.49340600
C       0.00000000      1.13336700      0.33992900
C      -0.00000000     -1.13336700      0.33992900
C       1.25441300     -0.78233900     -0.49340600
H       1.20544300      1.20489500     -1.50351600
H       2.16019200      1.17688700     -0.01817100
H       1.20544300     -1.20489500     -1.50351600
H       2.16019200     -1.17688700     -0.01817100
C       0.00000000      0.00000000      1.38887600
H      -0.89255700      0.00000000      2.02655300
H       0.89255700     -0.00000000      2.02655300
C      -1.25441300      0.78233900     -0.49340600
H      -2.16019200      1.17688700     -0.01817100
H      -1.20544300      1.20489500     -1.50351600
C      -1.25441300     -0.78233900     -0.49340600
H      -2.16019200     -1.17688700     -0.01817100
H      -1.20544300     -1.20489500     -1.50351600
H       0.00000000      2.15523000      0.73149700
H      -0.00000000     -2.15523000      0.73149700

#SMITH
schema = matrix.smith.input.v1
symmetrize = false
\end{lstlisting}

\begin{lstlisting}[basicstyle=\ttfamily\scriptsize,caption={Complete norbornane Gaussian ReadAllGIC input.}]
#p hf/sto-3g opt=(readallgic,calcfc,maxcycle=80) output=pickett

norbornane ESALC-I ReadAllGIC optimization

0 1
C       1.25441300      0.78233900     -0.49340600
C       0.00000000      1.13336700      0.33992900
C      -0.00000000     -1.13336700      0.33992900
C       1.25441300     -0.78233900     -0.49340600
H       1.20544300      1.20489500     -1.50351600
H       2.16019200      1.17688700     -0.01817100
H       1.20544300     -1.20489500     -1.50351600
H       2.16019200     -1.17688700     -0.01817100
C       0.00000000      0.00000000      1.38887600
H      -0.89255700      0.00000000      2.02655300
H       0.89255700     -0.00000000      2.02655300
C      -1.25441300      0.78233900     -0.49340600
H      -2.16019200      1.17688700     -0.01817100
H      -1.20544300      1.20489500     -1.50351600
C      -1.25441300     -0.78233900     -0.49340600
H      -2.16019200     -1.17688700     -0.01817100
H      -1.20544300     -1.20489500     -1.50351600
H       0.00000000      2.15523000      0.73149700
H      -0.00000000     -2.15523000      0.73149700

GIC001(Value=0.0) = 0.707106781187*(R(9,10))-0.707106781187*(R(9,11))
GIC002(Value=2.93156150623) = 0.707106781187*(R(9,10))+0.707106781187*(R(9,11))
GIC003(Value=1.09633572703) = R(1,6)
GIC004(Value=1.09633572703) = R(4,8)
GIC005(Value=1.09633572703) = R(12,13)
GIC006(Value=1.09633572703) = R(15,16)
GIC007(Value=1.09602638934) = R(4,7)
GIC008(Value=1.09602638934) = R(12,14)
GIC009(Value=1.09602638934) = R(15,17)
GIC010(Value=1.09602638934) = R(1,5)
GIC011(Value=1.09431690538) = R(2,18)
GIC012(Value=1.09431690538) = R(3,19)
GIC013(Value=1.564678) = R(1,4)
GIC014(Value=1.564678) = R(12,15)
GIC015(Value=0.0) = 0.707106781187*(R(3,4))-0.707106781187*(R(3,15))
GIC016(Value=0.0) = -0.707106781187*(R(1,2))+0.707106781187*(R(2,12))
GIC017(Value=4.13260235459) = 0.707106781187*(R(3,4))+0.707106781187*(R(3,15))
GIC018(Value=4.13260235459) = 0.707106781187*(R(1,2))+0.707106781187*(R(2,12))
GIC019(Value=0.0) = -0.707106781187*(R(2,9))+0.707106781187*(R(3,9))
GIC020(Value=4.12706024365) = 0.707106781187*(R(2,9))+0.707106781187*(R(3,9))
GIC021(Value=-0.0573424722848) = 0.816496580928*(A(5,1,6))-0.408248290464*(A(2,1,5))-0.408248290464*(A(2,1,6))
GIC022(Value=0.0136523036609) = 0.707106781187*(A(2,1,5))-0.707106781187*(A(2,1,6))
GIC023(Value=-0.0650891780287) = 0.816496580928*(A(5,1,6))-0.408248290464*(A(4,1,5))-0.408248290464*(A(4,1,6))
GIC024(Value=0.0195508295595) = 0.707106781187*(A(4,1,5))-0.707106781187*(A(4,1,6))
GIC025(Value=0.0337761539671) = 0.816496580928*(A(18,2,9))-0.408248290464*(A(18,2,12))-0.408248290464*(A(18,2,1))
GIC026(Value=0.0) = 0.707106781187*(A(18,2,12))-0.707106781187*(A(18,2,1))
GIC027(Value=0.0337761539671) = 0.816496580928*(A(19,3,9))-0.408248290464*(A(19,3,15))-0.408248290464*(A(19,3,4))
GIC028(Value=0.0) = 0.707106781187*(A(19,3,15))-0.707106781187*(A(19,3,4))
GIC029(Value=-0.0650891780287) = 0.816496580928*(A(7,4,8))-0.408248290464*(A(1,4,7))-0.408248290464*(A(1,4,8))
GIC030(Value=0.0195508295595) = 0.707106781187*(A(1,4,7))-0.707106781187*(A(1,4,8))
GIC031(Value=-0.0573424722848) = 0.816496580928*(A(7,4,8))-0.408248290464*(A(3,4,7))-0.408248290464*(A(3,4,8))
GIC032(Value=0.0136523036609) = 0.707106781187*(A(3,4,7))-0.707106781187*(A(3,4,8))
GIC033(Value=-0.0619062683331) = 0.816496580928*(A(10,9,11))-0.408248290464*(A(2,9,10))-0.408248290464*(A(2,9,11))
GIC034(Value=0.0) = 0.707106781187*(A(2,9,10))-0.707106781187*(A(2,9,11))
GIC035(Value=-0.0619062683331) = 0.816496580928*(A(10,9,11))-0.408248290464*(A(3,9,10))-0.408248290464*(A(3,9,11))
GIC036(Value=0.0) = 0.707106781187*(A(3,9,10))-0.707106781187*(A(3,9,11))
GIC037(Value=-0.0573424722848) = 0.816496580928*(A(13,12,14))-0.408248290464*(A(2,12,13))-0.408248290464*(A(2,12,14))
GIC038(Value=-0.0136523036609) = 0.707106781187*(A(2,12,13))-0.707106781187*(A(2,12,14))
GIC039(Value=-0.0650891780287) = 0.816496580928*(A(13,12,14))-0.408248290464*(A(15,12,13))-0.408248290464*(A(15,12,14))
GIC040(Value=-0.0195508295595) = 0.707106781187*(A(15,12,13))-0.707106781187*(A(15,12,14))
GIC041(Value=-0.0573424722848) = 0.816496580928*(A(16,15,17))-0.408248290464*(A(3,15,16))-0.408248290464*(A(3,15,17))
GIC042(Value=-0.0136523036609) = 0.707106781187*(A(3,15,16))-0.707106781187*(A(3,15,17))
GIC043(Value=-0.0650891780287) = 0.816496580928*(A(16,15,17))-0.408248290464*(A(12,15,16))-0.408248290464*(A(12,15,17))
GIC044(Value=-0.0195508295595) = 0.707106781187*(A(12,15,16))-0.707106781187*(A(12,15,17))
GIC045(Value=-0.0207227541011) = 0.632455532034*(A(3,4,1))-0.511667273602*(A(4,1,2))+0.195439507585*(A(1,2,9))+0.195439507585*(A(2,9,3))-0.511667273602*(A(9,3,4))
GIC046(Value=-0.063778079142) = 0.37174803446*(A(4,1,2))-0.601500955008*(A(1,2,9))+0.601500955008*(A(2,9,3))-0.37174803446*(A(9,3,4))
GIC047(Value=-0.0207227541011) = 0.632455532034*(A(3,15,12))-0.511667273602*(A(15,12,2))+0.195439507585*(A(12,2,9))+0.195439507585*(A(2,9,3))-0.511667273602*(A(9,3,15))
GIC048(Value=-0.063778079142) = 0.37174803446*(A(15,12,2))-0.601500955008*(A(12,2,9))+0.601500955008*(A(2,9,3))-0.37174803446*(A(9,3,15))
V0001(Value=-55.8772472164) = D(1,2,9,3)
V0002(Value=55.8772472164) = D(2,9,3,4)
V0003(Value=-35.2663208258) = D(9,3,4,1)
V0004(Value=0.0) = D(3,4,1,2)
V0005(Value=35.2663208258) = D(4,1,2,9)
V0006(Value=55.8772472164) = D(12,2,9,3)
V0007(Value=-55.8772472164) = D(2,9,3,15)
V0008(Value=35.2663208258) = D(9,3,15,12)
V0009(Value=0.0) = D(3,15,12,2)
V0010(Value=-35.2663208258) = D(15,12,2,9)
\end{lstlisting}

\subsection{Formic-acid--water fragment coordinate set}

The eight-atom complex has two fragments and closes at rank \(18/18\).  Six protected rows describe
the three relative translations and three relative orientations; the remaining
twelve rows span the intramolecular motions.  The resulting basis retains the intramolecular and fragment-motion families.

The orientation source follows the {translation--rotation
internal coordinate (TRIC)} construction of Wang and
Song:\cite{WangSong2016GeomeTRIC} least-squares superposition on a stored body reference defines the
optimal unit quaternion, its three-component exponential map defines the local
{coordinate set}, and analytic Cartesian derivatives define the {row of the Wilson B matrix}.  TRIC adds
six pose primitives to every subunit and delocalizes the augmented set.  The
coordinate set selects one reference fragment, retains
only relative pose variables, and protects their identities through rank and
projection under the declared molecular group.  The translation variables are
collective internal coordinates with analytic {rows of the Wilson B matrix}; no atomic Cartesian
variables are appended to the primitive pool in these weak-complex examples.
The body reference keeps the fragment orientations continuous during local optimization.

{A fixed-contact formic-acid--water structure compares pose,
contact--pose, and augmented-graph descriptions at common rank.  The two
hydrogen bonds close an auxiliary graph; the selected contact--pose coordinate set
retains finite fragment motion without adding valence torsions around that
cycle.}

\subsection{Closed double-hydrogen-bond diagnostic}

The nine-atom formamide--water geometry contains the contacts
N1--H3\(\cdots\)O7 and O7--H9\(\cdots\)O6.  Their H\(\cdots\)O distances are
2.0373 and 1.8876~\AA, and their donor--H--acceptor angles are
137.20\(^{\circ}\) and 150.26\(^{\circ}\), respectively.  The two contacts
increase the auxiliary cycle rank by one.  All bases were evaluated on this
same frozen Cartesian structure with analytic {rows of the Wilson B matrix} and without
symmetry projection (the complex is \(C_1\)).

The pose basis contains 15 intramolecular and six relative-fragment rows.
Table~\ref{tab:si-formamide-closed} reports the tangent-space
comparison for the closed-contact construction: two independent H-bond
distance rows were combined with four non-null pose directions after
projection against both the intramolecular and contact subspaces.  {The
augmented-graph control includes both hydrogen bonds in its valence adjacency.
The selected coordinate set uses finite pose functions for the residual directions.}
\(\kappa\) is computed after
normalizing every full {row of the Wilson B matrix}; \(\kappa_{\rm rel}\) uses the non-null
singular spectrum remaining after projection against the intramolecular
space.  Here \(f_B\) is the fraction of nonzero Wilson B matrix elements;
``atoms/row is the mean number of atoms in each coordinate. The projector
distance is the spectral norm of the difference from the pure-pose Cartesian
row-space projector.

\begin{table}[htbp]
\centering
\small
\caption{Analytic Wilson-matrix comparison for cyclic formamide--water.}
\label{tab:si-formamide-closed}
\begin{tabular}{@{}lrrrrrr@{}}
\toprule
Basis & Rank & \(\kappa\) & \(\kappa_{\rm rel}\) & \(f_B\) & atoms/row & projector distance \\
\midrule
Exponential-map pose & 21 & 8.052 & 2.050 & 0.503 & 4.52 & 0 \\
Two-contact tangent diagnostic & 21 & 7.964 & 2.778 & 0.429 & 3.86 & \(3.90\times10^{-15}\) \\
{Augmented-graph control} & 21 & 8.093 & 4.753 & 0.317 & 2.86 & \(3.05\times10^{-15}\) \\
\bottomrule
\end{tabular}
\end{table}

{All three descriptions close rank 21.  The contact--pose
diagnostic has a relative condition number of 2.778, compared with 2.050 for
pure pose and 4.753 for the augmented graph.}

\subsection{Primary multicenter bridges}

Closed noncovalent contacts and multicenter bonds have distinct chemical roles.  A recognized \(X\)--H--\(X'\) shared-hydrogen motif must be perceived
as a primary multicenter domain.  Proton
sharing between N/O/P/S centers and group-13 hydride bridges are separate
perception classes.  The latter includes homonuclear B--H--B and Al--H--Al and
heteronuclear Al--H--B.  Its source block contains two separately typed
center--H distances and the center--H--center bend.  Algorithmic recognition
is parameterized for qualified group-13 centers (B, Al, Ga, and In, with
heavier members admitted only when source data are available); raw distance equality and
center equivalence are not prerequisites.  These primary bridge tests use
qualified element-pair data, not the van-der-Waals screening reserved for
auxiliary contacts.  When several bridges share the same
pair of centers, a bridge-plane coupling completes their {remaining non-null directions of the Wilson B matrix}.
Symmetry projection forms exchange combinations only when supported by the
molecular group, and the global rank audit then operates on the complete
domain.  No bridge edge is passed through the auxiliary-contact cycle gate.

The validation set contains the \(C_{2v}\) shared-proton example (rank 21) and
\(D_{2h}\) diborane with two B--H--B domains (rank 18).  The diborane test
verifies all four bridge distances, full rank, assigned irreps, and the
point-group projector.  Additional homo- and heteronuclear group-13 examples
cover qualified B, Al, Ga, and In centers with element-pair-scaled distances;
they include analytic/finite-difference derivatives, atom-relabel covariance,
and deliberately unequal center--H distances.  Al--H--Al, Al--H--B, and the qualified group-13 centers are represented
by the same multicenter construction.



\subsection{\texorpdfstring{$\eta^3$}{eta3} allyl--palladium interaction-center coordinate set}

The ten-atom allyl--PdCl probe uses a frozen \(\eta^3\) center named
\texttt{C001} and spanning atoms 3--5.  The Pd--center distance is retained
as an independent coordinate; the coordinate set closes
at rank \(24/24\), with reference value 1.8000~\AA.  This checks the center-distance coordinate and its derivative.



These examples give Wilson B matrices of dimensions
\(18\times24\) and \(24\times30\), respectively.

\section{Domain-First Special-Coordinate Construction}

\subsection{Ownership order and conservative recognition}

The construction order tested in this work is
\[
\text{stretches}\longrightarrow\text{special-domain inventory}
\longrightarrow\text{complete special blocks}
\longrightarrow\text{ordinary residual coordinate set}.
\]
The special inventory contains ring and butterfly domains before local
\(XY_n\) domains, so an atom or primitive already owned by a cyclic source
space cannot be reclassified by a center rule.  \(XY_3\) requires a
four-coordinate acyclic center with three equivalent \(Y\) branches and one
distinct \(Z\) branch.  \(XY_2\) requires a three-coordinate center satisfying
the analogous (2+1) partition, and the center must not belong to a cycle.
Terminal atoms of the same element are exactly permutable.  Nonterminal
branches must additionally have identical rooted graph environments after
removing the center.  A fully symmetric \(XY_4\) or \(XY_3\) center has no
unique \(Z\) and is therefore not relabelled as \(XY_3-Z\) or \(XY_2-Z\),
respectively.

\subsection{Symmetry-resolved chemisorption block decomposition}

The illustrative main-text construction uses a central atom and the twelve
vertices of an ideal cuboctahedron for the fcc-like \(M_{13}\) substrate.  A
triangular face selects a \(C_{3v}\) site subgroup.  The symmetry-resolved
construction takes \(C_{3v}\) from the frozen state; the corresponding
\(C_1\) construction has the same block dimensions but no nontrivial irrep
decomposition.  The declared group remains fixed while the coordinate basis is
constructed.  In \(C_{3v}\) the
characters of the 39-dimensional substrate Cartesian representation are
\[
 \chi(E)=39,\qquad \chi(C_3)=0,\qquad \chi(\sigma_v)=3,
\]
and hence
\[
 \Gamma_{M_{13}}^{\mathrm{Cart}}=8A_1+5A_2+13E.
\]
Removing the polar translation representation \(A_1+E\) and the axial
rotation representation \(A_2+E\) gives the 33-dimensional substrate block
\[
 \Gamma_{M_{13}}^{\mathrm{Cart}/\mathrm{rigid}}=7A_1+4A_2+11E.
\]
The six natural internals of site-oriented NH\(_3\) span \(2A_1+2E\).
Relative translation and exponential-map orientation transform as \(A_1+E\)
and \(A_2+E\), respectively.  The complete typed direct sum is
therefore
\[
 (7A_1+4A_2+11E)\oplus(2A_1+2E)\oplus(A_1+A_2+2E)
 =10A_1+5A_2+15E,
\]
with dimension \(33+6+6=45=3(17)-6\).  The Cartesian substrate rows are
projected and derigidified before combination; they are not placed in a common
primitive pool with the molecular internals.

The same 33-dimensional substrate space can instead be generated without a
cluster bonding graph.  The complete \(M_{13}\) pair set contains 78 inverse
distances.  Its {Wilson B matrix} \(\Bmat_D\) has dimensions \(78\times39\) and
rank 33.  The source-space projector
\[
 \mathbf P_D=\Bmat_D\Bmat_D^+
\]
therefore has 33 unit eigenvalues and 45 null eigenvalues.  Retaining the
unit-eigenvalue eigenvectors and discarding the null space gives 33
non-redundant inverse-distance combinations.  A thin singular value
decomposition (SVD) of \(\Bmat_D\), a standard rank-revealing factorization,\cite{GolubReinsch1970}
or
deletion of the null eigenvalues of
\(\Bmat_D\Bmat_D^{\mathsf T}\), recovers the same subspace without explicitly
forming \(\mathbf P_D\); the latter additionally retains the nonzero metric
spectrum for conditioning diagnostics.  Symmetry projection gives the same
\(7A_1+4A_2+11E\) content as the derigidified Cartesian block.

Family-quota selection and the final Cartesian rank audit therefore see three
owned blocks and preserve the two-dimensional \(E\) subspaces for either
substrate realization.  This dimensionally closed construction is an
illustrative specification, not part of the present molecular validation
{set}.

\subsection{\texorpdfstring{\(XY_3\) and \(XY_2\)}{XY3 and XY2} coordinate blocks}

For \(X(Y_1Y_2Y_3)-Z\), the six primitive valence angles comprise three
\(ZXY_a\) and three \(Y_aXY_b\) rows.  The resulting five-dimensional
orthonormal local bending block contains two \(E\)-type combinations of the
\(ZXY_a\) orbit, the symmetric \(A_1\)-type combination of the
\(Y_aXY_b\) orbit, and two \(E\)-type combinations of that orbit.  The
missing symmetric angular direction is redundant with the bonded local frame.
The torsional coordinate is the common rotation
\[
q_{\mathrm{rot}}^{XY_3}=(3N_W)^{-1/2}
\sum_{w\in W(Z)}\sum_{a=1}^{3}D(w,Z,X,Y_a).
\]
Thus the two differential directions present in a three-branch torsional span
are carried by the local deformation block rather than duplicated as rotor
coordinates.

For an acyclic \(X(Y_1Y_2)-Z\) center, the owned angular block contains
\[
\frac{2A(Y_1,X,Y_2)-A(Z,X,Y_1)-A(Z,X,Y_2)}{\sqrt{6}},
\qquad
\frac{A(Z,X,Y_1)-A(Z,X,Y_2)}{\sqrt{2}},
\]
together with the oriented out-of-plane row of the same center.  Its collective
rotor is
\[
q_{\mathrm{rot}}^{XY_2}=(2N_W)^{-1/2}
\sum_{w\in W(Z)}\sum_{a=1}^{2}D(w,Z,X,Y_a).
\]
The expected periodicities are three and two for equivalent terminal
\(XY_3\) and \(XY_2\) groups.  Heterosubstitution, a cyclic \(XY_2\)
center, a linear branch or the absence of an admissible \(W\) anchor triggers
an ordinary single-dihedral alternative.

\subsection{Collective-rotor checks}

The Python and Fortran implementations give the same ordered primitive sets
and Wilson B matrix rows for the \(XY_3\) case. Further checks include a six-term normalized \(XY_2\) rotor for the
methylamine control, exclusive ownership of the \(XY_2\) out-of-plane row,
and the expected behavior of cyclic, heterosubstituted, and non-unique-\(Z\) motifs.
The ring examples retain their family assignments and vibrational rank.

\subsection{Independent and derived coordinates}

Let \(\mathbf p\) be the ordered primitive-source vector and let
\(\mathbf C\) contain the stored coefficients of the accepted generalized
coordinates.  The rank-carrying coordinates and their {Wilson B matrix} are
\begin{equation}
 \q=\mathbf C\mathbf p,
 \qquad
 \Bmat_{\q}=\mathbf C\Bmat_{\mathbf p}.
 \label{eq:si-stored-coordinate-definition}
\end{equation}
A derived coordinate \(\mathbf y=f(\q)\) is evaluated through the stored chain
rule,
\[
 \Bmat_{\mathbf y}
 =\frac{\partial f}{\partial\q}\Bmat_{\q}
 =\frac{\partial f}{\partial\q}\mathbf C\Bmat_{\mathbf p}.
\]
This covers inverse and exponential distances, trigonometric torsion
functions, proton-transfer combinations, and polar puckering amplitudes.

\subsection{{Coordinates requested by the user}}

{A \texttt{PROTECT} request names a chemically defined base
coordinate, such as a distance, torsion, ring-puckering or butterfly motion,
relative fragment displacement, or atom--center distance. Its analytic Wilson
B matrix row enters the appropriate chemical family before ordinary rank
completion. It is retained only if independent of the rows already accepted;
an incompatible excess of protected directions is reported rather than
silently replacing another family. A duplicate request is rejected, while a
user-defined row equivalent to an automatically generated one uses a single
coordinate with its origin recorded. Molecular symmetry is checked on the
complete orbit of a requested coordinate before projection.}

{An \texttt{OBSERVABLE} request instead defines a derived
quantity \(y=f(\q)\), such as an inverse distance, a trigonometric torsion
function, a proton-transfer combination, or a puckering amplitude. It is
evaluated by the chain rule above and adds no direction to the active
coordinate basis. These two requests thus distinguish a motion that must
participate in rank selection from a quantity evaluated from the selected
motions.}

{The corresponding request section uses named coordinate
functions. For example, the following records protect one O--H distance
and derive its inverse without increasing the rank:}
\begin{lstlisting}[basicstyle=\ttfamily\scriptsize]
SCHEMA smith.semantic.v1
PROTECT {
  oh = DISTANCE(1,2)
}
OBSERVABLE {
  inv_oh = INVERSE_DISTANCE(oh)
}
\end{lstlisting}
{These records belong to the optional
\texttt{SMITH\_SEMANTIC} section of the molecular-state input. Atom indices
refer to the ordered XYZ geometry; observables refer to named base
coordinates.}

\subsection{{Hessian transformation away from stationary points}}

{The Wilson B matrix determines the first-order transformation
of gradients. At a nonstationary geometry the exact Hessian transformation
also contains coordinate curvature. For independent coordinates \(q_i\), let
\(B_{ia}=\partial q_i/\partial x_a\),
\(D_{iab}=\partial^2q_i/(\partial x_a\partial x_b)\), and let
\(g_i^{(q)}\) and \(H_{ij}^{(q)}\) denote the internal-coordinate gradient
and Hessian. Then}
{
\begin{align}
g_a^{(x)} &= \sum_i B_{ia}g_i^{(q)}, \\
H_{ab}^{(x)} &= \sum_{ij}B_{ia}H_{ij}^{(q)}B_{jb}
                 + \sum_i g_i^{(q)}D_{iab}.
\end{align}
{With \(\mathbf C=\sum_i g_i^{(q)}\mathbf D_i\), the
corresponding Hessian on the full-rank internal tangent space is}
\[
\mathbf H^{(q)}=(\Bmat^+)^{\mathsf T}
                  (\mathbf H^{(x)}-\mathbf C)\Bmat^+ .
\]
}
{At an unconstrained stationary point the gradient vanishes
and so does \(\mathbf C\). Away from such a point, omitting this term would
attribute coordinate curvature to the potential-energy surface. Second
derivatives of the coordinate functions are evaluated only when a subsequent
analysis requires an off-stationary Hessian; they do not enter ESALC-I rank
selection or the stored first-derivative coordinate definition.}

{The implemented stationary-point transformation is
\(\mathbf H^{(q)}=(\Bmat^+)^{\mathsf T}\mathbf H^{(x)}\Bmat^+\).
It requires the Cartesian Hessian and the Wilson B matrix at that geometry;
no second coordinate derivatives are evaluated.}

\subsection{Typed exact-rank reduction and local stability}

Suppose \(\Bmat_A\) contains the already accepted rows and \(\Bmat_F\) the
next homogeneous family.  The Cartesian projector onto the accepted row space
and the residual family matrix are
\[
 \mathbf P_A=\Bmat_A^+\Bmat_A,
 \qquad
 \widetilde{\Bmat}_F=\mathbf D_F\Bmat_F(\mathbf I-\mathbf P_A),
\]
where \(\mathbf D_F\) normalizes nonzero source rows.  A rank-revealing
{orthogonal--triangular (QR) factorization} or
thin SVD of \(\widetilde{\Bmat}_F\) retains its non-null complement without
mixing coefficients from another family.  The accepted blocks are then
audited together against the nonlinear-molecule target \(3N-6\) (or \(3N-5\)
for a linear molecule).

For sequential selection, an ordered normalized row \(\widehat{\mathbf b}_i\)
is accepted when
\[
 \rho_i=\left\|(\mathbf I-\mathbf P_A)
              \widehat{\mathbf b}_i^{\mathsf T}\right\|_2
 >\tau_{\mathrm{rank}}.
\]
{For the reported construction runs,
\(\tau_{\mathrm{rank}}=10^{-7}\) for these row-normalized residuals.}
Absolute and relative singular-value thresholds, the rejected residual, and
the candidate order are stored.  Normalization prevents a change of length or
angle units from deciding rank, while the declared order retains the intended
chemical priority.

Under a state-preserving relabelling \(\pi\),
\[
 \Bmat(\pi_*\mathcal P)=\mathbf U_\pi\Bmat(\mathcal P)\mathbf Q_\pi^{\mathsf T},
\]
with \(\mathbf Q_\pi\) permuting Cartesian components and \(\mathbf U_\pi\)
permuting or signing oriented sources.  Canonical atom keys, pivot order, and
sign conventions therefore preserve the same coordinate space under relabelling. Within a fixed
molecular state, the {rows of the Wilson B matrix} are continuous functions of geometry.  If every
accepted/rejected residual, singular-value gap, template score, and symmetry
closure test has a nonzero margin from its threshold, continuity guarantees a
neighbourhood with the same accepted families and irreducible subspaces.  A
vanishing margin or a changed topology/equivalence record is reported as a
basin boundary that calls for reassessment.

\subsection{Ring source policy}

Ring construction begins from ordered elementary cycles and their
shared atoms and bonds.  Endocyclic angle sources generate the \(ring deformation\) family,
whereas \(n-3\) signed curvilinear rows span the local puckering family of an
isolated \(n\)-membered ring and retain finite Wilson derivatives at
planarity.  Alternative height and out-of-plane {coordinate sets} support operations
requiring a more global representation.  Cremer--Pople components and their
amplitudes and phases are evaluated as conformational descriptors without
increasing the active rank.

For fused, spiro, bridged, and cage systems, per-cycle candidates are assembled
with explicit butterfly or other inter-ring rows and reduced once in the
complete Cartesian tangent space.  Shared motion is retained once, so the
polycyclic dimension is established globally rather than by adding isolated
\(n_r-3\) counts.  Ring deformation, puckering, and collective inter-ring
motion retain separate family labels.  When a minimal {coordinate set} is not closed
under literal permutation of its labels, molecular symmetry is closed on the
physical Wilson-row span of the ring family before projection.

\subsection{Exact molecular symmetry and family quotas}

For each homogeneous family \(F\), the frozen operations induce a signed
representation \(\Gamma_F(g)\).  Closure is checked before the character
projector
\[
 \Pi_\alpha^{(F)}=
 \frac{\ell_\alpha}{|G|}\sum_{g\in G}
 \chi_\alpha(g)^*\Gamma_F(g)
\]
is applied.  Numerical ranges are extracted by rank-revealing factorization,
and a canonical pivot/sign gauge fixes the reported basis inside degenerate
spaces.  {For a multidimensional projected subspace, the projector
is symmetrized numerically, its columns are visited in the fixed source order,
independent columns are orthonormalized, and the first coefficient larger in
magnitude than \(10^{-12}\) is made positive.  This fixes a reproducible basis
without discarding a degenerate partner.}  For a multidimensional irrep, complete partner subspaces are retained;
one apparently convenient component cannot be selected alone.  Local
pseudosymmetry may order ligand or ring sources, but only this exact molecular
operation set assigns molecular irreducible representations.

{For nearly degenerate rows within one irreducible subspace,
individual coefficient vectors can change orientation along a geometry path
while the subspace remains the same. Overlaps between successive projected
subspaces identify this change; an orthogonal alignment can provide continuous
coordinate vectors when a subsequent analysis requires them.}

Independent choices from symmetry-projected families can be collectively rank
deficient even when each family has its requested dimension.  Selection preserves the assigned number of stretch, bend, torsion, ring, and
fragment directions while reaching global rank. Dependent combinations are
exchanged within their families; the construction reports any incompatible
family dimensions.

\subsection{Fixed coordinate definition and stored quantities}

{The construction retains the ordered primitive set
\(\mathbf p=(p_1,\ldots,p_m)^{\mathsf T}\), its analytic source Wilson B
matrix \(\Bmat_P\equiv\Bmat_{\mathbf p}\), the coefficient matrix \(\mathbf C\) of the
selected coordinates, and the molecular-symmetry operations. Equation
\ref{eq:si-stored-coordinate-definition} gives their values and derivatives
at a geometry within the assigned basin.}
{The stored definition also includes each primitive's family,
the selected and rejected candidates, their rank residuals, protected
coordinates, numerical tolerances, symmetry operations with atom
permutations, and the coefficients of local and molecular symmetry
combinations. These quantities identify the same coordinates whenever the
molecular state is evaluated again. Sparse analytic source rows are assembled
through Eq.~\ref{eq:si-stored-coordinate-definition}; a dense Wilson B matrix
is needed only for the rank and validation checks.}
{For SOOP and other geometry-dependent primitive functions,
\(\Bmat_P\) includes the derivatives of their defining coefficients; the
coordinate identities and the number of retained rows remain fixed. The
sparse-row density quoted in the validation tables is
\(\operatorname{nnz}(\Bmat)/(r\,3N)\), where \(r\) is the selected rank.}

{The molecular state supplied to this construction contains the
topology, rings, fragments, interaction centers, and symmetry operations.
When present, continuous atom and bond descriptors distinguish chemically
similar local environments: effective atomic numbers combine nuclear charge
with coordination, distance-derived bond order, angular rigidity, and
delocalization; synthon descriptors include charge, covalency, strain, and
separate first- and second-\(\pi\) contributions. These descriptors guide
ligand equivalence, ring orientation, and local torsion selection. Their
perception criteria and reassessment thresholds follow the published method\cite{Barone2026Perception}
and are implemented in the coordinate builder. The values used for construction remain part of
the fixed molecular state.}

\subsection{Properties of the ordered rank construction}

{Let \(\mathcal P=(p_1,\ldots,p_m)\) be the candidates in their
stored family and within-family order, and let \(r=3N-\dim T\), where
\(T\) is the space of overall translations and rotations. The target is
\(r=3N-6\) for a nonlinear molecule and \(r=3N-5\) for a linear molecule.
The following properties refer to a fixed molecular state with a complete
primitive row space.}

\paragraph{Rank completeness.}
{Write \(\mathcal R_k=\operatorname{span}\{\mathbf b(p_1),
\ldots,\mathbf b(p_k)\}\). At step \(k\), projection against the rows already
retained separates \(\mathbf b(p_k)\) into a component in
\(\mathcal R_{k-1}\) and an orthogonal residual. Retaining each nonzero
residual increases the accepted span by one dimension; rejecting a zero
residual leaves that span unchanged. Induction over the ordered candidates
therefore gives}
\begin{equation}
 {\operatorname{span}\Bmat_{\rm selected}
 =\operatorname{span}\Bmat_P,\qquad
 \operatorname{rank}\Bmat_{\rm selected}=\operatorname{rank}\Bmat_P.}
 \label{eq:si-rank-property}
\end{equation}
{The same conclusion holds for the familywise QR or singular-value
reduction in Sec.~S2.8: each retained residual subspace spans the contribution
of that family modulo the previously selected span. When
\(\operatorname{rank}\Bmat_P=r\), Eq.~\ref{eq:si-rank-property} yields exactly
\(r\) independent directions. A deficient primitive pool or incompatible
family dimensions is reported as such. The numerical acceptance test uses
row-normalized residuals and \(\tau_{\rm rank}=10^{-7}\), as specified above.}

\paragraph{Reproducibility for a fixed molecular state.}
{Primitive generation, family priority, atom ordering,
row normalization, rank thresholds, and symmetry operations are specified
before selection. Each subsequent choice is a deterministic function of
these quantities. A fixed input therefore yields the same selected coordinate
subspaces and coefficient convention, or the same reported rank or symmetry
incompatibility. This statement concerns the specified procedure; it does
not assert that a non-redundant basis is mathematically unique.}

\paragraph{Priority coordinates.}
{A protected ring, fragment, or interaction-center coordinate
is considered before ordinary rank-completion candidates. If its analytic
Wilson row is nonsingular and independent of the protected rows already
retained, it enters the final span. An ordinary stretch, bend, or torsion
considered later cannot remove it. If independent protected rows exceed
\(r\), the requested coordinate model is inconsistent and is reported rather
than reduced by discarding a protected motion. This follows directly from
the fixed order and the residual test.}

\paragraph{Atom relabelling.}
{Consider a permutation \(\pi\) that preserves elements,
geometry, topology, continuous descriptors, fragments, interaction centers,
symmetry operations, and ordering keys. Let \(\mathbf Q_\pi\) permute
Cartesian components and \(\mathbf U_\pi\) permute or sign oriented primitive
rows. The analytic derivatives obey}
\begin{equation}
 {\Bmat_P(\pi\mathbf x)=
 \mathbf U_\pi\Bmat_P(\mathbf x)\mathbf Q_\pi^{\mathsf T}.}
 \label{eq:si-relabel-property}
\end{equation}
{Both matrices are orthogonal, so row norms, residual ranks,
and the dimensions of projected symmetry subspaces are preserved. The same
family and candidate decisions follow in the corresponding order. Thus the
accepted coordinate spaces and symmetry blocks transform with the atom
relabeling; their printed signs are fixed by the stored coefficient convention.}

\paragraph{Scaling within a family.}
{Multiplying every row in one homogeneous source block
\(\Bmat_F\) by the same nonzero factor \(\lambda\) leaves its normalized
row directions unchanged up to sign. It also multiplies
\(\Bmat_F\Bmat_F^{\mathsf T}\) by \(\lambda^2\), without changing its
eigenspaces. Consequently the retained subspace is unchanged when the
singular-row and rank thresholds are not crossed. The fixed family order
keeps this operation within one chemically defined class.}

\paragraph{Local stability.}
{Suppose topology, atom equivalence, fragment membership,
symmetry operations, family order, and protected candidates remain fixed
within a structural neighbourhood of radius \(\delta_{\rm str}\) around
\(\mathbf x_0\). Let \(f_a(\mathbf x)\) denote the finite set of numerical
decision quantities, including row norms, rank residuals, singular values,
template scores, and coefficient-sign choices. For a threshold \(\tau_a\),
assume a positive margin
\(\mu_a=|f_a(\mathbf x_0)-\tau_a|>0\) and a local Lipschitz bound
\(|f_a(\mathbf x)-f_a(\mathbf y)|\leq L_a\|\mathbf x-\mathbf y\|\).
Every displacement satisfying}
\begin{equation}
 {\|\mathbf x-\mathbf x_0\|<\delta,
 \qquad 0<\delta<\min\!\left(
 \delta_{\rm str},\min_{a:L_a>0}\frac{\mu_a}{2L_a}\right)}
 \label{eq:si-local-stability}
\end{equation}
{leaves each decision on the same side of its threshold.
Therefore the primitive identities, retained dimensions, family labels, and
irreducible subspaces stay fixed, while their analytic Wilson rows vary
continuously. A zero margin or a changed molecular state calls for reassessment.
For repeated or nearly degenerate symmetry
subspaces, continuity concerns the subspace; individual printed coefficient
vectors may rotate within it, as discussed above.}

\subsection{Deterministic symmetry and derivative evaluation}

{Local symmetry-adapted linear combinations are formed from
chemically equivalent primitive orbits. The totally symmetric sum is stored
first, followed by orthogonal difference combinations in a fixed atom order.
Ring operations must preserve the ring's atom and bond descriptors; a shared
edge in a fused system is assigned once before the complete ring source space
is reduced. Ordinary stretches, bends, and exocyclic torsions then fill the
remaining rank in their declared order. For exocyclic bonds the representative
dihedral is chosen from the continuous ligand priorities and local
connectivity; only equivalent signed rows with stable derivatives are combined.}

{For an ordered orbit of \(m\) equivalent sources
\(p_1,\ldots,p_m\), the stored local combinations are}
{
\begin{align}
 q_{A_1}&=\frac{1}{\sqrt m}\sum_{i=1}^{m}p_i,\\
 q_{k+1}&=\frac{\sum_{i=1}^{k}p_i-kp_{k+1}}{\sqrt{k(k+1)}},
 \qquad k=1,\ldots,m-1.
 \label{eq:si-local-combinations}
\end{align}
}
{The difference rows are orthonormal, have zero coefficient
sum, and remain within one primitive family. Ligands are first grouped by
their continuous chemical descriptors and radial shell. The nominal
equivalence tolerances used in the local-pseudosymmetry tests are
\(5\times10^{-4}\) for effective atomic number and \(10^{-3}\)~\AA\ for
center--ligand distance; explicitly relaxed controls are identified in
Table~\ref{tab:si-named-coordination-controls}.}

{For a high-coordination template \(T\), let \(\mathbf c\)
and \(\mathbf t_T\) be the sorted ligand-pair cosine vectors of length
\(M=n(n-1)/2\). The rotationally and permutation-invariant score is}
\begin{equation}
 {s_T=\left[\frac{1}{M}\sum_{k=1}^{M}
 (c_k-t_{T,k})^2\right]^{1/2}.}
 \label{eq:si-template-score}
\end{equation}
{A named template is selected when its score is at most 0.12
and the gap to the next score is at least 0.02. Other geometries use the
ligand-equivalence construction with the same final rank test. The score,
gap, selected coordinate family, and thresholds accompany the fixed
coordinate definition.}

{The final molecular-symmetry operation acts within a closed
coordinate family. Its character projector selects the complete subspace of
each irreducible representation, including every partner of a degenerate
representation. Projected columns are visited in source order, orthonormalized,
and assigned a reproducible sign. The recorded operation matrices and atom
permutations permit the same projector to be evaluated at later geometries
in the basin. For a geometry path, overlaps between consecutive projected
subspaces separate a harmless rotation of coefficient vectors from a change
in the molecular assignment.}
{Because the operations form a representation on a closed
family, character orthogonality makes projectors for distinct irreducible
representations mutually orthogonal and idempotent. Their ranges sum to the
family's accepted row space. Keeping the full nonzero range of each projector
therefore preserves both the family rank and complete symmetry partners;
the final Cartesian rank audit detects dependencies between families.}

{Analytic primitive Wilson rows are evaluated at the current
geometry and combined with the stored coefficients through
Eq.~\ref{eq:si-stored-coordinate-definition}. A derived quantity
\(y=f(\mathbf q)\) uses
\(\partial y/\partial\mathbf x=(\partial f/\partial\mathbf q)\Bmat\),
so it does not change rank or family ownership. The first derivative of the
Wilson B matrix is evaluated only by a calculation that requires a Hessian
away from a stationary point, using the curvature term of Sec.~S2.7.}

\section{Gaussian Generalized Internal Coordinate (GIC) Input}

\subsection{\texttt{ReadAllGIC} input}

Gaussian 16 Rev. C.02 reads the complete ESALC-I set through
\texttt{ReadAllGIC}; native calculations use its built-in redundant
coordinates. The supplied block contains the independent coordinates,
needed ring and center functions, and frozen symmetry partners. The
MP2/6-31G(d)\cite{HariharanPople1973Polarization} comparison uses \texttt{CalcHFFC,Tight} in both coordinate
paths. The \texttt{External} keyword provides energies and gradients for
the comparisons with a common external potential.

\subsection{Scaling of coordinates for Gaussian}

Wilson-metric normalization is used during rank selection and molecular
projection because it gives a well-conditioned, scale-neutral independence
test.  Gaussian's optimizer initializes forces and force constants on the
conventional scale of its primitive internal coordinates.  The exported {coordinate set}
therefore applies one diagonal, optimizer-specific rescaling.  If final row
\(k\) has primitive coefficients \(c_{kp}\), the Gaussian coordinate is
\begin{equation}
 q_k^{(G)}=s_k q_k,
 \qquad
 s_k=\left(\sum_p c_{kp}^2\right)^{-1/2}.
 \label{eq:gaussian-optimizer-scaling}
\end{equation}
Thus every exported primitive-coefficient vector has unit Euclidean norm.
The transformation is nonsingular and diagonal, so it preserves the Wilson-row
space, exact rank, sparsity pattern, family ownership, and irrep assignment.
It only presents the same ESALC-I {coordinate set} in the scale expected by the Gaussian
force and initial-Hessian model.  This scaling is applied uniformly to every exported ESALC-I coordinate set.

This point is directly visible for the open-graph water dimer.  Its normalized
H-bond row is \(2^{-1/2}R_{\mathrm{H\cdots O}}\); Eq.~\ref{eq:gaussian-optimizer-scaling}
exports it as \(R_{\mathrm{H\cdots O}}\).  The resulting initial force constant
is 0.04663, equal to the native distance scale, while the complete {coordinate set}
retains rank 12 and its \(C_s\) irrep partition.  The same single rescaling rule
is applied to every ESALC-I input in the controlled comparison below.

\subsection{Curvilinear ring coordinates in Gaussian}

The signed local out-of-plane compatibility {coordinate set} uses

\[
U(j;a,b,c)=\arcsin\!\left[
\frac{\hat r_{ja}\cdot(\hat r_{jb}\times\hat r_{jc})}
{|\hat r_{jb}\times\hat r_{jc}|}\right],
\]

with the stored atom order fixing the sign. Gaussian expresses this
function through
\(U(I,J,K,L)\mapsto D(J,I,L,K)\).  The generated signed puckering rows are
written through algebraically equivalent ordered curvilinear kernels.  The coordinate retains its family, coefficients, rank, and symmetry.  When molecular symmetry is retained, all non-totally
symmetric ESALC-I rows are written with \texttt{Frozen}; for norbornane this
leaves 15 active \(A_1\) rows and 36 frozen rows in the full 51-coordinate
basis.

\section{Construction Times}

\subsection{Reference-path construction-scaling data}

{The polycyclic aromatic hydrocarbon (PAH) series is included in}
Table~\ref{tab:si-construction-cost}, {which gives} the numerical data underlying the
construction-cost summary in the main text.  Times are medians of five single-process Python
runs on a Linux workstation with two AMD EPYC 7543 32-core processors and
1.0~TiB {random-access memory (RAM)}.  Build time includes construction, rank reduction, optional
symmetry orientation, and output; the \(\Bmat\) column is one analytic
evaluation of the stored sparse coordinate set.  Peak Python allocation was measured
with \href{https://docs.python.org/3/library/tracemalloc.html}{\texttt{tracemalloc}}.  These measurements document
{benchmark breadth}, sparse evaluation, and the reference algorithm.

For the medium-to-large {Parametrized Model 6 (PM6)} comparison
summarized in the main text, the
tetraphenylporphyrin, docetaxel, and paclitaxel ESALC-I paths required 12.31,
79.99, and 55.79~s, respectively, for 17, 116, and 72 optimization steps.
These elapsed times were measured in a single process on an Apple~M3~Pro
(arm64, 36~GB memory, macOS~26.5.2), using the same machine for the
construction audit.  The timings include the complete Gaussian--PM6 route and
are provided as hardware-specific execution data; the scientific comparison
is based on the common PM6 level and the reported optimization paths.

\begin{table}[htbp]
\centering
\scriptsize
\setlength{\tabcolsep}{3pt}
\caption{Single-process scaling of the coordinate construction.}
\label{tab:si-construction-cost}
\begin{tabular}{@{}llrrrrrr@{}}
\toprule
Series & System & \(N\) & Candidates & Rank & Build / s & \(\Bmat\) / s & Peak MiB \\
\midrule
PAH & Benzene & 12 & 30 & 30 & 0.131 & 0.0054 & 0.23 \\
PAH & Naphthalene & 18 & 48 & 48 & 0.139 & 0.0093 & 0.28 \\
PAH & Anthracene & 24 & 66 & 66 & 0.243 & 0.0139 & 0.51 \\
PAH & Phenanthrene & 24 & 66 & 66 & 0.211 & 0.0131 & 0.52 \\
PAH & Pyrene & 26 & 78 & 72 & 0.256 & 0.0162 & 0.59 \\
PAH & Coronene & 36 & 120 & 102 & 0.722 & 0.0343 & 1.08 \\
Ring/cage & Cubane & 16 & 54 & 42 & 0.400 & 0.0139 & 0.45 \\
Ring/cage & Norbornane & 19 & 52 & 51 & 0.099 & 0.0120 & 0.29 \\
Ring/cage & Cyclooctane & 24 & 66 & 66 & 0.142 & 0.0201 & 0.53 \\
Ring/cage & Camphor & 27 & 76 & 75 & 0.145 & 0.0151 & 0.44 \\
Ring/cage & Spiro & 31 & 87 & 87 & 0.253 & 0.0225 & 0.75 \\
Special & Ferrocene \(D_{5h}\) & 21 & 651 & 57 & 3.324 & 0.0642 & 2.83 \\
Special & Ferrocene \(D_{5d}\) & 21 & 656 & 57 & 3.182 & 0.0653 & 2.59 \\
\bottomrule
\end{tabular}
\end{table}

{
\section{Construction and Ring-Coordinate Validation}

The {validation set} underlying the representative
main-text table is collected in Table~\ref{tab:si-validation-full}.
Horizontal rules separate the individual systems so that rank closure,
coordinate family and numerical audit can be followed independently.
Here SALC denotes a symmetry-adapted linear combination, \(RPck\) denotes
a ring-puckering coordinate, and \(\Delta_\mathrm{PF}\) is the
Python--Fortran difference in the numerical check. The point group (PG)
and finite-difference (FD) errors follow the definitions above.

\scriptsize
\setlength{\tabcolsep}{3pt}
\begin{longtable}{@{}>{\raggedright\arraybackslash}p{0.14\linewidth}>{\raggedright\arraybackslash}p{0.24\linewidth}>{\raggedright\arraybackslash}p{0.07\linewidth}>{\raggedright\arraybackslash}p{0.15\linewidth}>{\raggedright\arraybackslash}p{0.31\linewidth}@{}}
\caption{Rank, symmetry, and derivative checks for the ESALC-I construction.}\label{tab:si-validation-full}\\
\toprule
System & Stress test & PG & Rank/closure & Numerical check \\
\midrule
\endfirsthead
\toprule
System & Stress test & PG & Rank/closure & Numerical check \\
\midrule
\endhead
Benzene & Aromatic ring projector & \(D_{6h}\) & 30/30; 3 \(RPck\) & \(\Delta_\mathrm{PF}=6.67\times10^{-9}\); \(\epsilon_\mathrm{SALC}=2.22\times10^{-16}\) \\
\midrule
Pyrrole & Heteroaromatic ring projector & \(C_{2v}\) & 24/24; 2 \(RPck\) & \(\Delta_\mathrm{PF}=6.92\times10^{-9}\); \(\epsilon_\mathrm{SALC}=2.22\times10^{-16}\) \\
\midrule
Azulene & Non-benzenoid fused ring, butterfly & \(C_{2v}\) & 48/48; 6 \(RPck\) & \(\Delta_\mathrm{PF}=6.44\times10^{-9}\); \(\epsilon_\mathrm{SALC}=2.22\times10^{-16}\) \\
\midrule
Norbornane & Bridged saturated ring & \(C_{2v}\) & 51/51; 3 \(RPck\) & \(\Delta_\mathrm{PF}=6.22\times10^{-9}\); \(\epsilon_\mathrm{SALC}=2.22\times10^{-16}\) \\
\midrule
Cubane & High-symmetry cage & \(O_h\) & 42/42; cyclic bends & \(\Delta_\mathrm{PF}=6.41\times10^{-9}\); \(\epsilon_\mathrm{SALC}=2.22\times10^{-16}\) \\
\midrule
Ferrocene & \(D_{5h}\) metallocene projector & \(D_{5h}\) & 57/57 & \(\Delta_\mathrm{PF}=6.71\times10^{-9}\); \(\epsilon_\mathrm{SALC}=3.33\times10^{-16}\) \\
\midrule
Spiro system & Spiro ring source spaces & \(D_2\) & 87/87; 6 \(RPck\) & \(\Delta_\mathrm{PF}=6.97\times10^{-9}\); \(\epsilon_\mathrm{SALC}=2.22\times10^{-16}\) \\
\midrule
Cyclooctane & Large-ring puckering & \(D_2\) & 66/66; 5 \(RPck\) & \(\Delta_\mathrm{PF}=9.84\times10^{-9}\); \(\epsilon_\mathrm{SALC}=2.22\times10^{-16}\) \\
\midrule
Pagodane & Cage and fused-ring source spaces & \(D_{2h}\) & 114/114 & \(\epsilon_\mathrm{FD}=4.02\times10^{-10}\); stretches and cyclic bends covered \\
\midrule
Indane & Fused aromatic/alicyclic rings & \(C_s\) & 51/51 & \(\epsilon_\mathrm{FD}=4.77\times10^{-10}\); butterfly and \(RPck\) rows covered \\
\midrule
Saccharin & Planar fused heterocycle, butterfly--puckering phase & \(C_s\) & 45/45; 5 \(RPck\) & \(\epsilon_\mathrm{FD}=3.91\times10^{-10}\); full post-projector \(\Bmat\) rank 45 \\
\midrule
Formic acid--water & Weak-complex fragment model & \(C_1\) & 18/18; 6 protected & three translations plus three orientations; direct \(\Bmat\)-rank \(=18\) \\
\midrule
Water dimer & Hydrogen-bonded fragment model & \(C_1\) & 12/12; 6 protected & standalone Cartesian input; three translations plus three orientations \\
\midrule
\(\eta^3\) allyl--PdCl & Supplied \(\eta^3\) ligand center & \(C_1\) & 24/24; 1 protected & \(q_{\eta^3}=1.8000\)~\AA; analytic center--atom row retained \\
\midrule
Open weak-complex contacts & Intermolecular and intramolecular pseudobond contacts & \(C_s\), \(C_{3v}\), \(C_1\) & rank closed & formic acid--water, water dimer, Li\(_2\)--Ar, borane--ammonia and glycine I/II retain contact origin \\
\midrule
Formamide--water & Closed double-H-bond diagnostic & \(C_1\) & 21/21 for all three bases & \(\kappa_{\rm rel}=2.050\), 2.778, and 4.753 for pose, contact--pose, and {augmented-graph bases} \\
\midrule
Shared-H bridge domains & Primary multicenter hyperedges & \(C_{2v}\), \(D_{2h}\) & 21/21 and 18/18 & central-proton and diborane tests require complete bridge distances, assigned irreps and {full rank of the Wilson B matrix} \\
\midrule
\bottomrule
\end{longtable}
\normalsize

Three molecular structures illustrate the ring-coordinate construction.  The
first is pagodane, \(\mathrm{C}_{20}\mathrm{H}_{20}\): structural perception identifies 50 covalent bonds, 11 rings, no aromatic atoms and a
\(D_{2h}\) reference assignment.  The primitive pool and ESALC-I basis both have Cartesian rank 114.  The second is indane,
\(\mathrm{C}_{9}\mathrm{H}_{10}\): structural perception identifies 20 bonds, two fused
rings, six aromatic atoms and a \(C_s\) reference assignment, while ESALC-I closes
at rank 51.

The third example is planar saccharin, retained specifically as a fused-ring
phase and global-rank validation.  Its \(C_s\) coordinate set contains one butterfly
and five ring-puckering coordinates, all in \(A''\), and closes at rank 45/45.
The butterfly coordinate is
\begin{equation}
 q_{\mathrm{bf}}=2^{-1/2}\!\left[D(6,1,2,13)-D(16,1,2,3)\right],
\end{equation}
whose relative minus sign supplies the inter-ring direction missing from the
five \(RPck\) rows.  Replacing it by the corresponding sum lowers the rank of
these six rows from six to five.  The validation therefore checks the complete
post-symmetry Cartesian \(\Bmat\), rather than accepting six individually
well-formed but collectively dependent expressions.

Central finite differences of every final ESALC-I row give maximum absolute
errors of \(4.02\times10^{-10}\), \(4.77\times10^{-10}\), and
\(3.91\times10^{-10}\) for pagodane, indane, and saccharin, respectively.  The
audit follows the complete chain rule and includes stretch, bend, cyclic-bend,
butterfly, and puckering families.

For every one of the 114 pagodane and 51 indane coordinates, the Cartesian
fingerprint is
\[
 \Delta\mathbf{x}_k=\mathbf{B}_0^+\mathbf{e}_k\Delta q_k,
 \qquad \mathbf{x}_{\pm}=\mathbf{x}_0\pm\Delta\mathbf{x}_k .
\]
An optional minimum-mass-norm form uses
\(\mathbf{M}^{-1/2}(\mathbf{B}_0\mathbf{M}^{-1/2})^+\); its origin records
average, default-isotope, Hessian, or named-isotopologue masses.  Step reduction
preserves the frozen topology and orientation, limits the largest displacement
to 0.10~\AA, and limits each reference-bond change to 15\%.  These symmetric
frames freeze all other ESALC-I coordinates only to first order and are not nonlinear
internal-to-Cartesian realizations.

Near-planar fused rings are covered by signed puckering and butterfly rows whose
first derivatives remain finite at the planar reference.  Symmetry acts on
their physical Wilson-row spaces within the corresponding ring families, and
the complete \(\Bmat\) rank is then checked across all families.  Indane tests
the derivative accuracy of this construction; saccharin tests the coupled case
in which butterfly and puckering rows share the same molecular irrep.  Together
they verify continuity, intrafamily projection, and global rank closure.

The diagnostic record contains the right-inverse metric, mass
origin, normalized-row conditioning, Cartesian amplification, participation
ratio, significantly moving atoms, dominant atom, local domain and pseudogroup,
independent local/global total-symmetry labels, canonical ring order, primitive
coefficients, and topology-check count.  Display amplitudes are independently
normalized for viewing.

\subsection{Ring Puckering Branch}

{Figure~\ref{fig:si-cyclobutane-pucker} shows the fixed-frame
cyclobutane scan generated from one retained ring-puckering coordinate.}
\begin{figure}[htbp]
\centering
\includegraphics[width=0.9\linewidth]{figures/sonic_puckering_scan.pdf}
\caption{{Cyclobutane puckering energy along the fixed-frame
ESALC-I coordinate at the HF/STO-3G level.}}
\label{fig:si-cyclobutane-pucker}
\end{figure}


{The cyclobutane \(D_{2d}\) puckering profile at HF/STO-3G has a
minimum near \(q=0.04\)~\AA{} and rises smoothly on either side.}
A {discrete-variable-representation (DVR) or potential-energy-surface
(PES)} model may replace this base coordinate by a Cremer--Pople
amplitude or phase while retaining the frozen analytic chain rule.


\subsection{Portable Chair-to-Boat Puckering Export}

The cyclohexane ring-flip calculation tests a generated puckering path.
From graph and geometry, ESALC-I generates a fixed-frame path from the chair
\(q_3\) component to the boat \(q_2\) component.  Each point is exported both
as ordinary Cartesians and as an automatically generated \texttt{ReadAllGIC}
coordinate set, then evaluated with the {Becke three-parameter
Lee--Yang--Parr (B3LYP)} functional and 6-31+G* basis.

All five coordinate sets have rank \(48/48\).  Their Cartesian and GIC energies differ
by at most \(3.0\times10^{-3}\,\mu E_h\), or about
\(8\times10^{-6}\)~kJ~mol\(^{-1}\). The energy rise reflects the
unrelaxed fixed-frame deformation.

Figure~\ref{fig:si-cyclohexane-export} compares the five-point energy
traces and their differences. The two representations follow the same
fixed-frame deformation.

\begin{figure}[htbp]
\centering
\includegraphics[width=0.92\linewidth]{figures/cyclohexane_puckering_equivalence.pdf}
\caption{{Cyclohexane chair-to-boat fixed-frame scan at
B3LYP/6-31+G*. Cartesian and generated ReadAllGIC inputs give equivalent
energies; the lower panel shows their difference.}}
\label{fig:si-cyclohexane-export}
\end{figure}

{Independent PBE0/def2-SVP single-point calculations with
ORCA\cite{Neese2020ORCA} and Molpro\cite{Werner2020Molpro} on three of these
geometries give the relative energies in
Table~\ref{tab:si-cyclohexane-independent}. The unrelaxed path is a test of
portable geometries, not a chair-to-boat barrier.}

\begin{table}[htbp]
\centering
\caption{{Cyclohexane fixed-frame energies from two independent programs at PBE0/def2-SVP.}}
\label{tab:si-cyclohexane-independent}
\begin{tabular}{crrrr}
\toprule
Path point & ORCA $E/E_h$ & ORCA $\Delta E$ & Molpro $E/E_h$ & Molpro $\Delta E$ \\
\midrule
$p_{00}$ ($\lambda=0.00$) & $-235.422463982$ & 0.000 & $-235.422142275$ & 0.000 \\
$p_{02}$ ($\lambda=0.50$) & $-235.321345943$ & 265.485 & $-235.321038853$ & 265.447 \\
$p_{04}$ ($\lambda=1.00$) & $-235.103704314$ & 836.903 & $-235.103335256$ & 837.028 \\
\bottomrule
\end{tabular}
\par\smallskip\footnotesize $\Delta E$ is in kJ~mol$^{-1}$ relative to $p_{00}$.
\end{table}



\clearpage
\section{Stationary-Hessian Comparisons}

\subsection{Stationary-Hessian locality control}

The heat maps in the main text use the optimized cyclobutane HF/STO-3G
stationary point; the
corresponding Gaussian route is \texttt{HF/STO-3G Opt Freq NoSymm}.  The same Cartesian
Hessian is used for every representation.  Six external directions are removed by an
SVD to form the 30-dimensional orthonormal Cartesian vibrational control.  For
{the delocalized control, the Gaussian coordinate definitions supply
72 native primitives: 12 distances, 24 angles, and 36 dihedrals.  Their
Wilson B matrix is evaluated at the stationary geometry in bohr and
radians using the analytic derivatives in geomeTRIC~1.1.1.\cite{WangSong2016GeomeTRIC}
The 30 eigenvectors of \(\Bmat\Bmat^{\mathsf T}\) with positive
eigenvalues above \(10^{-10}\) of the largest form the delocalized basis.
{The recovered molecular group is \(D_{2d}\); its eight
operations match the stationary geometry within \(1.7\times10^{-8}\)~\AA.
Projection of
the 30 delocalized directions gives six \(A_1\), two \(A_2\), three \(B_1\),
five \(B_2\), and fourteen \(E\) partner directions.  The seven nearly
degenerate eigenvalue pairs are the \(E\) subspaces; rotating their partners
preserves molecular symmetry while changing individual force-matrix elements.
The displayed matrix orders these directions by \(D_{2d}\) irrep.}}  For
each ESALC-I construction, the frozen analytic {Wilson B matrix} is evaluated at
the stationary Cartesian geometry and defines the
Cartesian realization
\[
 \mathbf A=\Bmat^{\mathsf T}
 (\Bmat\Bmat^{\mathsf T})^+,
 \qquad
\mathbf F_q=\mathbf A^{\mathsf T}\mathbf F_x\mathbf A.
\]
This frame match is essential: the internal reference is stored in a canonical
principal-axis orientation, whereas the Cartesian Hessian remains in the
Gaussian Cartesian frame.  Combining the latter with a {Wilson B matrix} evaluated
in the former would be an invalid rotation mismatch.  In the matched frame the
raw \texttt{NoSymm} Hessian is already invariant under the recovered \(D_{2d}\)
operations to a relative Frobenius residual below \(7.1\times10^{-9}\), so no
group averaging is applied.
The optimized reference has negligible forces, so this stationary-point
transformation requires only the stored first-order {Wilson B matrix}.  Table
\ref{tab:hessian-locality} reports the off-diagonal statistics plotted in the
main text; diagonal elements are excluded.
{The second and fourth maps group delocalized and
symmetry-adapted ESALC-I rows by \(D_{2d}\) irrep.  The fourth also groups
rows by coordinate family (stretch, bend, cyclic bend, and ring pucker);
solid and dotted separators mark the respective boundaries.}

\begin{table}[htbp]
\centering
\caption{Normalized off-diagonal Hessian coupling for cyclobutane.  All
representations have dimension and rank 30.}
\label{tab:hessian-locality}
\begin{tabular}{lccc}
\toprule
Representation & Mean & 95th percentile & Maximum \\
\midrule
Cartesian vibrational & 0.15099 & 0.39380 & 0.59591 \\
{Delocalized Gaussian primitives} & {0.03475} & {0.21239} & {0.95506} \\
ESALC-I coordinate set & 0.02794 & 0.10497 & 0.74501 \\
Symmetry-adapted ESALC-I & 0.04253 & 0.31296 & 0.93174 \\
\bottomrule
\end{tabular}
\end{table}

For the symmetry-oriented basis, the 120 allowed intra-irrep entries have
mean 0.15419, 95th percentile 0.51606, and maximum 0.93174.  The remaining 315
inter-irrep entries have maximum unrounded magnitude
\(4.6\times10^{-8}\) and therefore plot as exact zero at six-decimal storage
precision.  Symmetry orientation thus produces disjoint blocks while retaining the
chosen coordinate basis within each block.

Isolated chemically meaningful couplings remain in the maximum. The mean and
percentile quantify overall coupling, and symmetry gives the visible block
structure.
{Primitive internal functions give the most nearly diagonal
chemically local force-constant representation, but are redundant.  The
delocalized control removes redundancy by combining all primitive families.
ESALC-I keeps the families distinct through local SALCs and ring-deformation
coordinates before exact-rank selection and final molecular-symmetry
orientation.  Its measured coupling is low while its rows retain chemical
assignments.  The 72 primitive rows span a rank-30 space.}

\subsection{Norbornane stationary-Hessian comparison}

{The norbornane geometry was optimized at the
Hartree--Fock/STO-3G level with PySCF~2.14.0, starting from the geometry in
Listing~S1.  The final energy is \(-268.8938238457\) hartree; the largest
Cartesian gradient component is \(6.2\times10^{-6}\) hartree/bohr.  The
lowest eigenvalue of the Cartesian Hessian projected onto the 51-dimensional
vibrational space is positive (0.00238 hartree/bohr\(^2\)).  One analytic
Cartesian Hessian is used for every panel of main-text Fig.~10.}

{The redundant primitive pool contains 20 distances, 42
angles, and 72 dihedrals.  Its Wilson B matrix has rank 51.  Delocalized
directions are the eigenvectors of \(\Bmat\Bmat^{\mathsf T}\) with eigenvalues
above \(10^{-10}\) of the largest, as in the cyclobutane control.  The local
and symmetry-adapted ESALC-I sets each contain 51 independent coordinates.
The delocalized eigenvectors are symmetry-pure under the recovered
\(C_{2v}\) operations to numerical precision (minimum irrep-projector
fraction greater than 0.999999).
The Cartesian, delocalized, and ESALC-I force matrices use the same
transformation defined above.  The optimized geometry has \(C_{2v}\)
symmetry; the final ESALC-I set contains 15 \(A_1\), 11 \(A_2\), 12 \(B_1\), and
13 \(B_2\) directions.  Table~\ref{tab:norbornane-hessian} gives the coupling
statistics.  The largest inter-irrep coupling in the final set is
\(3.4\times10^{-8}\).}

\begin{table}[htbp]
\centering
\caption{Norbornane normalized off-diagonal Hessian coupling.  All
non-redundant representations have rank 51.}
\label{tab:norbornane-hessian}
\begin{tabular}{lccc}
\toprule
Representation & Mean & 95th percentile & Maximum \\
\midrule
Cartesian vibrational & 0.10733 & 0.33673 & 0.55853 \\
Delocalized internals & 0.04042 & 0.25869 & 0.82044 \\
Local ESALC-I & 0.02424 & 0.09951 & 0.71758 \\
Symmetry-adapted ESALC-I & 0.03456 & 0.25471 & 0.74491 \\
\bottomrule
\end{tabular}
\end{table}

\section{Local-Pseudosymmetry Extended Validation}

The accepted template score is denoted \(s\), its separation from the
second-best template by \(\Delta s\), the acceptance ceiling by
\(s_{\max}=0.12\), and the minimum separation by
\(\Delta s_{\min}=0.02\).  {The point group (PG) and template
scores} for all 20 ideal/distorted models appear in Table~\ref{tab:si-local-full}.
{Coordination-number (CN) 5--9 models} were tested.  Python and the separate Fortran implementation agree on every local
group, place the local totally symmetric row first, and close the target rank.

The shared domain matcher covers all coordination numbers above three.  The
CN~4 entries are tetrahedral ($T_d$) and square-planar ($D_{4h}$); the
CN~5--9 entries are the trigonal-bipyramidal/square-pyramidal, octahedral/
trigonal-prismatic, pentagonal-bipyramidal/capped-octahedral,
square-antiprismatic/dodecahedral, and tricapped-prismatic/capped-square-
antiprismatic alternatives listed in Table~\ref{tab:si-local-full}.  CN~10 is covered by bicapped
square-antiprismatic and bicapped-dodecahedral templates, CN~11 by capped-
icosahedral and octadecahedral templates, and CN~12 by icosahedral,
cuboctahedral, and hexagonal-bipyramidal templates.  CN~13 and CN~14 are
covered by capped-cuboctahedral and bicapped-hexagonal-antiprismatic templates,
respectively.  Unmatched environments
use the available generic rank-closing construction.  These assignments are
made before coordinate export and do not depend on Gaussian.

{For the capped-octahedral coordination-7 model, the cap lies on a local
threefold axis. The three vertices of the capped face and the three opposite
vertices are separate ligand orbits. Applying the triplet SALCs shown in the
main text to each orbit gives
\(3A_1+2E\) stretching directions; the complete eight-atom model has rank
18, as reported in Table~\ref{tab:si-local-full}. The seven stretches leave
11 independent angular directions, selected from the valence-angle orbits
by the same rank test. The explicit angular SALCs in the main text have
local symmetry \(2A_1+A_2+4E\). At the ideal geometry the three
opposite-ligand angles are linear, so their ordinary angle derivatives are
singular; the remaining 18 angle sources span all 11 angular directions.
No dihedral is intrinsic to this bare coordination sphere; a bonded
substituent or accepted collective rotor supplies a distinct torsional
family when present.}

\begin{landscape}
\scriptsize
\begin{longtable}{@{}ccp{0.27\linewidth}ccrrrr@{}}
\caption{Complete high-coordination local-template validation.}\label{tab:si-local-full}\\
\toprule
CN & Local group & Template & Geometry & \(s\) & \(\Delta s\) & \(s_{\max}\) & \(\Delta s_{\min}\) & Rank \\
\midrule
\endfirsthead
\toprule
CN & Local group & Template & Geometry & \(s\) & \(\Delta s\) & \(s_{\max}\) & \(\Delta s_{\min}\) & Rank \\
\midrule
\endhead
5 & \(D_{3h}\) & trigonal bipyramidal & ideal & \(2.48\times10^{-17}\) & 0.217962 & 0.12 & 0.02 & 12 \\
5 & \(D_{3h}\) & trigonal bipyramidal & distorted & 0.007152 & 0.206418 & 0.12 & 0.02 & 12 \\
5 & \(C_{4v}\) & square pyramidal & ideal & \(5.30\times10^{-17}\) & 0.217962 & 0.12 & 0.02 & 12 \\
5 & \(C_{4v}\) & square pyramidal & distorted & 0.012670 & 0.202571 & 0.12 & 0.02 & 12 \\
6 & \(O_h\) & octahedral & ideal & 0 & 0.375781 & 0.12 & 0.02 & 15 \\
6 & \(O_h\) & octahedral & distorted & 0.006421 & 0.364851 & 0.12 & 0.02 & 15 \\
6 & \(D_{3h}\) & trigonal prismatic & ideal & \(9.72\times10^{-17}\) & 0.375781 & 0.12 & 0.02 & 15 \\
6 & \(D_{3h}\) & trigonal prismatic & distorted & 0.005769 & 0.368504 & 0.12 & 0.02 & 15 \\
7 & \(D_{5h}\) & pentagonal bipyramidal & ideal & 0 & 0.174704 & 0.12 & 0.02 & 18 \\
7 & \(D_{5h}\) & pentagonal bipyramidal & distorted & 0.005878 & 0.165714 & 0.12 & 0.02 & 18 \\
7 & \(C_{3v}\) & capped octahedral & ideal & \(5.93\times10^{-17}\) & 0.174704 & 0.12 & 0.02 & 18 \\
7 & \(C_{3v}\) & capped octahedral & distorted & 0.005464 & 0.167037 & 0.12 & 0.02 & 18 \\
8 & \(D_{4d}\) & square antiprismatic & ideal & \(3.85\times10^{-17}\) & 0.296745 & 0.12 & 0.02 & 21 \\
8 & \(D_{4d}\) & square antiprismatic & distorted & 0.004965 & 0.290042 & 0.12 & 0.02 & 21 \\
8 & \(D_{2d}\) & dodecahedral-like & ideal & \(1.33\times10^{-16}\) & 0.296745 & 0.12 & 0.02 & 21 \\
8 & \(D_{2d}\) & dodecahedral-like & distorted & 0.008716 & 0.284759 & 0.12 & 0.02 & 21 \\
9 & \(D_{3h}\) & tricapped trigonal prismatic & ideal & \(2.62\times10^{-17}\) & 0.112154 & 0.12 & 0.02 & 24 \\
9 & \(D_{3h}\) & tricapped trigonal prismatic & distorted & 0.004833 & 0.105571 & 0.12 & 0.02 & 24 \\
9 & \(C_{4v}\) & capped square antiprismatic & ideal & 0 & 0.112154 & 0.12 & 0.02 & 24 \\
9 & \(C_{4v}\) & capped square antiprismatic & distorted & 0.004434 & 0.106386 & 0.12 & 0.02 & 24 \\
\bottomrule
\end{longtable}
\end{landscape}
\normalsize

{Table~\ref{tab:si-named-coordination-controls} gives the two
named high-coordination controls and their thresholds.}  Distorted
\(\mathrm{SF_6}\) requires deliberately relaxed effective-atomic-number and
distance thresholds, demonstrating that pseudosymmetry is an explicit user
choice rather than an implicit relabelling.

\begin{table}[htbp]
\centering
\scriptsize
\caption{Named high-coordination controls and their declared thresholds.}
\label{tab:si-named-coordination-controls}
\begin{tabular}{@{}lccrrrrr@{}}
\toprule
System & Template & Group & \(s\) & \(\Delta s\) & \(\tau_Z\) & \(\tau_r\)/\textup{\AA} & Rank \\
\midrule
Distorted \(\mathrm{SF_6}\) & octahedral & \(O_h\) & 0.006472 & 0.364380 & 0.1 & 0.05 & 15/15 \\
\([\mathrm{ZrF_8}]^{4-}\) & square antiprismatic & \(D_{4d}\) & 0.000017 & 0.296729 & 0.0005 & 0.001 & 21/21 \\
\bottomrule
\end{tabular}
\end{table}

The tests also include CN 10--14 generic alternatives and an ambiguous near-tie control.

\subsection{GOPPEG all-ESALC-C optimization}

{GOPPEG is the Cambridge Structural Database (CSD) identifier for an 81-atom zirconium organometallic complex in the transition metal quantum mechanics (tmQM) dataset.\cite{BalcellsSkjelstad2020tmQM} Its distorted starting geometry uses a complete
\(C_1\) ESALC-C basis of rank \(237/237\) after removal of overall
translation and rotation. The LINK
optimizer uses second-generation geometry, frequency, and noncovalent
extended tight binding (GFN2-xTB) energies and analytic
gradients.\cite{BannwarthEhlertGrimme2019GFN2xTB} Table~\ref{tab:goppeg-sycart}
reports the converged result. Mixed ESALC-I/ESALC-C handling is implemented
but is not used in this work.}

\begin{table}[htbp]
\centering
\caption{GOPPEG optimization with a complete ESALC-C basis.}
\label{tab:goppeg-sycart}
\begin{tabular}{lr}
\toprule
Quantity & Result \\
\midrule
Optimization steps & 32 \\
Energy--gradient evaluations & 47 \\
Maximum Cartesian gradient / \(E_h/a_0\) & \(1.8498\times10^{-4}\) \\
Root-mean-square Cartesian gradient / \(E_h/a_0\) & \(4.1180\times10^{-5}\) \\
\bottomrule
\end{tabular}
\end{table}

\section{Coordinate Motifs}

{Table~\ref{tab:si-motif-tests} collects the coordinate-motif
checks; FD denotes finite-difference derivatives of the Wilson B matrix.}

\begin{table}[htbp]
\centering
\scriptsize
\setlength{\tabcolsep}{3pt}
\caption{Numerical checks for representative coordinate families.}
\label{tab:si-motif-tests}
\begin{tabularx}{\linewidth}{@{}>{\raggedright\arraybackslash}p{0.20\linewidth}>{\raggedright\arraybackslash}p{0.30\linewidth}>{\raggedright\arraybackslash}X@{}}
\toprule
Molecular motif & ESALC-I coordinates & Numerical result \\
\midrule
Formamide transfer & \(R_{\rm NH}\), \(R_{\rm OH}\), \(\Sigma r\), \(\Delta r\) & For \(R_{\rm NH}=1.02\) and \(R_{\rm OH}=1.21\)~\AA, \(\Sigma r=2.23\) and \(\Delta r=-0.19\)~\AA; sparse rows equal the exact sum/difference. \\
Li\(_2\)--Ar scan & Two protected translation axes, polar radius and elliptic quadratic & At \((x,y)=(3,4)\), \(R=5.000000\), \((x/7.084)^2+(y/4.718)^2=0.898137\), and the elliptic row has six nonzero Cartesian entries. \\
\(XY_3\) local block & Projected homogeneous \(C_{3v}\) bend block & Rank 3/3; ordered rows are \(A_1\), \(E\), \(E\), with exactly one totally symmetric row. \\
\(XY_2\) local block & Complete acyclic \((2+1)\) branch domain & Methylamine retains two in-plane rows, one oriented out-of-plane row, and one six-term normalized collective rotor with exclusive family ownership. \\
Weak-complex rotor & Protected fragment translations and orientations & Water-dimer coordinate set rank 12/12; {Wilson B matrix} shape \(12\times18\); finite rigid-step residual below 0.02 in the six protected variables. \\
Ferrocene cyclopentadienyl torsion & Ring-center/Fe coordinates over \(0,9,18,27,36^\circ\) & Rank and analytic-\(\Bmat\) rank are 57 at all five points; two ring centers persist; sparse density is 0.095--0.101 and the smallest singular value remains above \(10^{-8}\). \\
Cyclohexane ring flip & Three symmetry-adapted ring-puckering directions & Rank 48/48; three puckering rows and three frozen Gaussian expressions. \\
Reduced path subgroup & Same cyclohexane source rebuilt with retained group \(C_1\) & Rank remains 48/48; global high-symmetry labels are not imposed along the path and scan torsions remain uncombined. \\
\bottomrule
\end{tabularx}
\end{table}

These checks establish the automatic construction and analytic rows that feed
the subsequent optimizer tests.

\section{Optimization Results}

\subsection{Common tight-binding starting geometries for noncovalent systems}

The noncovalent comparisons start from complexes optimized with the {second-generation geometries,
frequencies, and noncovalent extended tight-binding method (GFN2-xTB)}
(\href{https://github.com/grimme-lab/xtb}{xTB} 6.7.1), using the standard
geometry-optimization driver and tight optimization thresholds.  The method is
described in Ref.~\cite{BannwarthEhlertGrimme2019GFN2xTB}.  {Table~\ref{tab:gfn2xtb-starting-geometries}
records the resulting iteration counts and energies.}

\begin{table}[htbp]
\centering
\scriptsize
\caption{GFN2-xTB preprocessing of the noncovalent starting geometries.}
\label{tab:gfn2xtb-starting-geometries}
\begin{tabular}{lrr}
\toprule
System & GFN2-xTB iterations & Final energy / $E_h$ \\
\midrule
Water dimer & 39 & $-10.148095861237$ \\
Ammonia--water & 43 & $-9.505231322616$ \\
Ammonia dimer & 9 & $-8.855304465487$ \\
Formamide--water & 8 & $-15.729368709572$ \\
Formic acid--water & 11 & $-16.369999724051$ \\
Imidazole--water, acceptor & 35 & $-19.376267789565$ \\
Imidazole--water, donor & 34 & $-19.377553067244$ \\
Uracil--water, NH donor & 57 & $-29.701535010796$ \\
Uracil--water, carbonyl acceptor & 77 & $-29.696547792107$ \\
Benzene--water & 17 & $-20.950381937980$ \\
Benzene dimer & 36 & $-31.765815629075$ \\
Formamide cyclic dimer & 14 & $-21.313639924814$ \\
\bottomrule
\end{tabular}
\end{table}

{Table~\ref{tab:noncovalent-sonic-summary} records the ESALC-I
optimization outcomes and observed basin behavior for the same systems.}

\begin{table}[htbp]
\centering
\scriptsize
\setlength{\tabcolsep}{3.5pt}
\caption{Final ESALC-I noncovalent runs.  The basin column records the observed
topological behavior during optimization.}
\label{tab:noncovalent-sonic-summary}
\begin{tabularx}{\linewidth}{@{}l l r l X@{}}
\toprule
System & {Coordinate family} & Steps & Outcome & Basin behavior \\
\midrule
Water dimer & Pseudobond & 8 & Converged & H-bond 2--4 retained \\
Ammonia--water & Pseudobond & 9 & Converged & H-bond 3--4 retained (1.877\(\rightarrow\)1.911~\AA) \\
Ammonia dimer & Quaternion & 7 & Converged & Equivalent-H exchange only \\
Formic acid--water & Pseudobond & 11 & Converged & H-bond 5--6 retained (1.737\(\rightarrow\)1.663~\AA) \\
Formamide--water & Pseudobond & 16 & Converged & Contact 6--9 retained \\
Benzene--water & Quaternion & 8 & Converged & OH--\(\pi\) basin retained \\
Benzene dimer & Quaternion & 12 & Converged & Equivalent ring-atom exchange only \\
Formamide cyclic dimer & Quaternion & 7 & Converged & Cyclic motif retained \\
Imidazole--water, acceptor & Pseudobond & 17 & Converged & H-bond 5--10 retained (1.968\(\rightarrow\)1.845~\AA) \\
Imidazole--water, donor & Pseudobond & 20 & Converged & H-bond 9--11 retained (1.880\(\rightarrow\)1.930~\AA) \\
Uracil--water, NH donor & Pseudobond & 12 & Converged & Contact 9--13 retained \\
Uracil--water, carbonyl acceptor & Pseudobond & 19 & Converged & H-bond 12--15 retained (1.886\(\rightarrow\)1.896~\AA) \\
\bottomrule
\end{tabularx}
\end{table}

The benzene dimer supplies a dispersion-bound test case.  The
cyclic formamide dimer was initialized from the literature \(C_{2h}\) double
hydrogen-bonded motif and then optimized independently with GFN2-xTB; its
optimized geometry is used in the comparison.
For benzene--water, the optimized GFN2-xTB structure was rigidly rotated so
that the water hydrogens face the ring before the three optimizer calculations;
{this geometry is the common starting structure.}
The optimized Cartesian geometries are
the starting structures for the subsequent ESALC-I, Gaussian redundant, and
PySCF/geomeTRIC comparisons.\cite{WangSong2016GeomeTRIC}

The construction tests above establish rank and symmetry.  The external tests
then test direct use of the complete ESALC-I set in Gaussian 16 Rev. C.02
and compare optimization routes on a common potential-energy surface.  The
paired Gaussian diagnostic in
Table~\ref{tab:gaussian-matched-coordinate-paths} compares automatic ESALC-I
generation with Gaussian's built-in redundant-coordinate generation inside
the Gaussian optimization framework.  The controlled PySCF comparison below
instead holds the
potential and starting structures fixed across PySCF/geomeTRIC,
Gaussian/ESALC-I, and Gaussian/native routes.  These tests define
their stated finite protocols and span both the covalent and noncovalent
optimization methods used in this work.

\subsection{Gaussian 16 Rev. C.02 \texttt{ReadAllGIC} interface}

ESALC-I generates the complete \texttt{ReadAllGIC} block for
Gaussian 16 Rev. C.02. The \texttt{External} keyword connects this
optimization machinery to arbitrary energy and gradient sources.
The coordinate-consumption probes in
Table~\ref{tab:readallgic-opt}: curvilinear ring rows without a corresponding
commercial GIC name are written through algebraically equivalent ordered
four-point kernels, and unsupported long multi-dihedral expressions are
expanded into the accepted syntax.  The final geometry is re-evaluated against
the generated coordinate definitions (Sec.~S3).  Every ESALC-I row remains
explicit, with non-totally symmetric rows marked \texttt{Frozen} whenever
symmetry is imposed.  At the {Hartree--Fock (HF)} level with the
Slater-type-orbital three-Gaussian (STO-3G) minimal basis, the three compatibility probes reach
stationary points in three to nine steps.  {The Gaussian
\(NTRed/NRank\) fields report the redundant-coordinate count and rank when
available.}

\begin{table}[htbp]
\centering
\caption{{Gaussian 16 Rev. C.02 ESALC-I optimizations at
HF/STO-3G.}}
\label{tab:readallgic-opt}
\begin{tabular}{lccccc}
\toprule
System & PG & \ESALCI\ rank & Gaussian \(NTRed/NRank\) & Steps & \(E_\mathrm{final}\) / \(E_h\) \\
\midrule
Cyclohexane & \(D_{3d}\) & 48/48 & -- & 5 & \(-231.482671171\) \\
Norbornane & \(C_{2v}\) & 51/51 & -- & 9 & \(-268.893823819\) \\
Cubane & \(O_h\) & 42/42 & -- & 3 & \(-303.781399977\) \\
\bottomrule
\end{tabular}
\end{table}

A symmetry-enabled norbornane variant retains the full rank-closing block but
activates only its 15 \(A_1\) rows.  It reaches a \(C_{2v}\)-preserving stationary
point in eight steps (Secs.~S1 and S3).

Additional small and medium-size checks used the same starting geometries for
the native and ESALC-I paths.  The settings were MP2/6-31G(d), \texttt{Tight}
for malonaldehyde and diborane, and HF/STO-3G without \texttt{Tight} for the
two ferrocene structures; all four calculations used standard guess Hessians.
All eight jobs terminated normally; the malonaldehyde geometry relaxed from
the declared \(C_{2v}\) structure to \(C_s\) in both paths.
{Table~\ref{tab:additional-gaussian-checks} gives the paired
step counts and endpoint energy differences.}

\begin{table}[htbp]
\centering
\caption{Supplementary paired Gaussian checks beyond the main comparison.}
\label{tab:additional-gaussian-checks}
\begin{tabular}{lccccl}
\toprule
System & Method & ESALC-I steps & Native steps & Final PG & $|\Delta E|/E_h$ \\
\midrule
Malonaldehyde ($C_{2v}$) & MP2/6-31G(d), Tight & 10 & 10 & $C_s$ & $2.2\times10^{-7}$ \\
Diborane ($D_{2h}$) & MP2/6-31G(d), Tight & 5 & 5 & $D_{2h}$ & $7.0\times10^{-8}$ \\
Ferrocene ($D_{5h}$) & HF/STO-3G & 6 & 6 & $D_{5h}$ & $2.6\times10^{-7}$ \\
Ferrocene ($D_{5d}$) & HF/STO-3G & 6 & 6 & $D_{5d}$ & $2.8\times10^{-7}$ \\
\bottomrule
\end{tabular}
\end{table}

\subsection{Optimization comparison on a common potential-energy surface}

{Five systems use the Perdew--Burke--Ernzerhof hybrid (PBE0)
functional\cite{AdamoBarone1999PBE0} with the def2-SVP basis\cite{WeigendAhlrichs2005Def2} and third-generation dispersion
(D3) with Becke--Johnson damping\cite{GrimmeEhrlichGoerigk2011D3BJ} and the
Axilrod--Teller--Muto three-body term [D3(BJ)-ATM]:}
glycine, the \(C_s\) malonaldehyde minimum, the \(C_{2v}\) malonaldehyde
{transition state (TS)}, the water dimer, and formamide--water.  Every route starts from
the same symmetrized GFN2-xTB geometry.  PySCF~2.14.0 supplies {restricted
Kohn--Sham (RKS)} energies and
analytic gradients, including the D3(BJ)-ATM term, either directly to
geomeTRIC~1.1.1 or through Gaussian's \texttt{External} interface.

The PySCF/geomeTRIC calculation uses its default TRIC path and standard initial
Hessian.  The Gaussian 16 Rev. C.02 ESALC-I route uses the generated
\texttt{ReadAllGIC} coordinates; the native route uses built-in redundant internals.  Neither Gaussian route uses
\texttt{CalcFC}.  All three routes use root-mean-square (RMS) and maximum gradient thresholds of
\(3.0\times10^{-4}\) and \(4.5\times10^{-4}\,E_h a_0^{-1}\), and RMS and
maximum displacement thresholds of \(1.2\times10^{-3}\) and
\(1.8\times10^{-3}\,a_0\).  The displacement thresholds are converted to
\(6.3501\times10^{-4}\) and \(9.5252\times10^{-4}\)~\AA{} for geomeTRIC,
which also applies its \(10^{-6}\,E_h\) energy-change criterion.  Gaussian
applies its corresponding force/displacement termination test.  The use of
each optimizer's standard Hessian model makes this an end-to-end route
comparison.

The coordinate protocol is fixed before optimization.  Covalent systems use
the automatic natural-internal {coordinate set}; the open water-dimer contact graph uses
pseudobonds; and the closed formamide--water graph uses the exponential-map
relative-pose block.  Every ESALC-I row is rescaled according to
Eq.~\ref{eq:gaussian-optimizer-scaling}.  Non-totally symmetric rows are retained
as \texttt{Frozen}, while the native Gaussian calculations use Gaussian's
molecular-symmetry machinery.
{Table~\ref{tab:pyscf-cross-optimizer} compares the resulting
optimization paths from the common starting geometries.}

\begin{table}[htbp]
\centering
\scriptsize
\setlength{\tabcolsep}{3.5pt}
\caption{Energy--gradient evaluations for five common-potential optimizations.}
\label{tab:pyscf-cross-optimizer}
\begin{tabularx}{\linewidth}{@{}lXrrrr@{}}
\toprule
System & Contact type & ESALC-I & Gaussian native & Established optimizer & Rank/point group \\
\midrule
Glycine & covalent & 4 & 4 & 6 & 24/24, \(C_s\) \\
Malonaldehyde minimum & covalent & 7 & 7 & 10 & 21/21, \(C_s\) \\
Malonaldehyde TS & covalent & 5 & 6 & 9 & 21/21, \(C_{2v}\) \\
Water dimer & open/pseudobond & 10 & 10 & 12 & 12/12, \(C_s\) \\
Formamide--water & closed/exponential map & 11 & 16 & 15 & 21/21, \(C_1\) \\
\midrule
Total evaluations & & 37 & 43 & 52 & \\
\bottomrule
\end{tabularx}
\end{table}

All calculations terminate normally in the intended basin.  The minimum
ESALC-I metric eigenvalues remain positive throughout; the smallest over the
five paths is 0.0162.  The largest aligned ESALC-I--native endpoint
{root-mean-square deviation (RMSD)} is
0.0018~\AA{}, and the largest energy difference among any two route endpoints
is \(11.1\,\mu E_h\).  The three covalent ESALC-I paths use 16 energy--gradient evaluations,
and the two weak-complex paths use 21. Table~\ref{tab:pyscf-cross-optimizer}
gives the counts for the established optimization routes. All five systems
start from common geometries and use the same potential-energy surface.

\subsection{Native outcomes for the boundary tests}

The native-redundant Al$_{13}$ and Al$_{13}^{+}$ jobs both reach normal
termination and satisfy the stationary-point criteria after 45 and 53 steps,
respectively.  In both cases the optimization leaves the initial \(O_h\)
subspace and terminates in \(C_1\), at
\(-3106.67435751\,E_h\) and \(-3106.75576280\,E_h\).  The corresponding
symmetry-oriented inverse-distance calculations terminate after six and five
steps while retaining \(O_h\).  The paths end in different symmetry subspaces, illustrating the role of
explicit symmetry orientation.

A native-redundant pyrene--NH$_3$ optimization from an
unrelaxed assembly likewise reaches normal termination after 34 steps, using
Gaussian's negligible-force criterion.  At that point the nitrogen atom lies
approximately 25.7~\AA{} from the pyrene plane and the final energy is
\(-659.780742043\,E_h\): the calculation has reached the dissociated
\(C_s\) limit.  The fragment-pose calculation below starts from separately relaxed fragments.

\subsubsection{\texorpdfstring{Preoptimized-fragment adsorption protocol: pyrene--NH$_3$}{Preoptimized-fragment adsorption protocol: pyrene--NH3}}

To distinguish relaxation within each fragment from relaxation of the
interfragment pose, pyrene and NH$_3$ were optimized separately at the cheaper HF/STO-3G
with their automatically generated ESALC-I coordinates.  The isolated jobs
required four and five Gaussian geometry steps, respectively.  Gaussian can
place independently optimized fragments in different standard Cartesian
frames; therefore each final fragment geometry was rigidly aligned to its own
starting frame before the complex was assembled.  The assembled 84-coordinate
coordinate set was re-evaluated at the combined geometry and supplied unchanged to a
full ESALC-I optimization.  The latter reached normal termination in 13 steps,
with \(NTRed/NRank=84/84\), final point group \(C_s\), and
\(E=-659.781527448\,E_h\).  Rigid alignment preserves each optimized fragment geometry.

\subsection{Paired native-redundant and ESALC-I coordinate paths}

Gaussian 16 Rev. C.02 reads the generated ESALC-I set through
\texttt{ReadAllGIC} and builds redundant coordinates for the native path.  For each
system, the native-redundant and ESALC-I jobs use the same
starting Cartesian geometry, MP2/\allowbreak6-31G(d),
\texttt{CalcHFFC},\texttt{Tight}, Gaussian
symmetry, and geometry-optimization protocol.  In the ESALC-I path Gaussian receives the complete \(3N-6\) generated block: all totally symmetric
rows are active and every non-totally symmetric row is also supplied,
explicitly marked \texttt{Frozen}.  All eighteen jobs reached normal termination.

The eight systems that remain in the same symmetry require 49 aggregate steps
with native redundant internals and 46 with ESALC-I.
Final energies for those eight pairs differ by at most
\(2.1\times10^{-7}\,E_h\).  Four cases tie; ferrocene requires one fewer ESALC-I
step, the sumanene minimum one more, alanine one fewer, and 2-deoxyribose two fewer.
The paired paths have similar endpoint energies and convergence counts.

\begin{table}[htbp]
\centering
\scriptsize
\setlength{\tabcolsep}{3.0pt}
\caption{{Paired Gaussian MP2/6-31G(d) coordinate paths with
\texttt{CalcHFFC,Tight} and molecular symmetry.}}
\label{tab:gaussian-matched-coordinate-paths}
\begin{tabularx}{\linewidth}{@{}lccccXX@{}}
\toprule
System & ESALC-I rank & Active/Frozen & ESALC-I steps & Native steps &
ESALC-I PG path & Native PG path \\
\midrule
Cubane & 42/42 & 2/40 & 3 & 3 & \(O_h\) & \(O_h\) \\
Saccharin & 45/45 & 30/15 & 6 & 6 & \(C_s\) & \(C_s\) \\
Ferrocene & 57/57 & 4/53 & 6 & 7 & \(D_{5h}\) & \(D_{5h}\) \\
Pagodane & 114/114 & 18/96 & 7 & 7 & \(D_{2h}\) & \(D_{2h}\) \\
Sumanene minimum & 93/93 & 31/62 & 7 & 6 & \(C_3\) & \(C_3\) \\
Sumanene planar & 93/93 & 11/82 & 7 & 14 &
\(D_{3h}\) & \(D_{3h}\!\rightarrow C_1\!\rightarrow C_{3v}\) \\
C$_{60}$ & 174/174 & 2/172 & 5 & 5 & \(I_h\) & \(I_h\) \\
Alanine I & 33/33 & 33/0 & 7 & 8 & \(C_1\) & \(C_1\) \\
2-Deoxyribose & 51/51 & 51/0 & 5 & 7 & \(C_1\) & \(C_1\) \\
\bottomrule
\end{tabularx}
\end{table}

{The planar-sumanene paths reach different stationary subspaces
and are reported individually.}  The native-redundant job
starts at \(D_{3h}\), is recognized as \(C_1\) from steps 2--7, and reaches
the \(C_{3v}\) bowl after 14 steps.  ESALC-I remains in \(D_{3h}\) through
normal termination at step~7.  The row demonstrates explicit symmetry control
in addition to the iteration comparison.  The native final energy is, as
expected, essentially that of the independently started bowl minimum.

In exact arithmetic, a gradient evaluated at a symmetric geometry is totally
symmetric.  {Gaussian's built-in redundant GICs are not explicitly}
organized into irreducible-representation blocks.  Finite-precision coordinate
evaluation, linear solves, Hessian updates, and back-transformation can
therefore introduce small non-totally symmetric components, and active
nonsymmetric directions allow them to accumulate.  ESALC-I assigns every row to
an irreducible representation, activates only the totally symmetric block, and
retains every other rank-closing row as \texttt{Frozen}; numerical components
outside the declared symmetry subspace are consequently projected out at every
step.  ESALC-I preserves the requested point group in all nine cases.

Cubane and pagodane test
high-symmetry cage and many-cycle construction, saccharin tests a fused
heterocycle, and ferrocene tests automatically perceived ring-center/haptic and
fragment-frame families.  Sumanene adds a seven-ring fused polycycle in its
\(C_3\) bowl and planar \(D_{3h}\) forms.  In both forms, 93 automatically
generated rows are reduced to exact rank and symmetry oriented before
writing.  The planar structure is a \(D_{3h}\) stationary-point
reference with one imaginary mode; the present calculation tests preservation
of that symmetry subspace.

The C$_{60}$ case starts from an ideal \(I_h\) truncated icosahedron.  Its
31-cycle basis generates 74 stretches, 46 cyclic bends, and 54 butterfly rows.
The resulting 174 coordinates retain \(NTRed/NRank=174/174\) and the full
120-operation \(I_h\) group through normal termination at step~5, with only
two \(A_g\) rows active and all other rows \texttt{Frozen}.  Pagodane likewise
retains rank 114/114 and \(D_{2h}\) through all seven steps.

The frozen ferrocene fragment translation and orientation rows are written
by their exact analytic reference {rows of the Wilson B matrix}.  This is locally exact for a
frozen coordinate and avoids a singular angle--axis helper at coincident
reference frames.  All 57 ESALC-I rows remain present, of which 53 are frozen.

The additional covalent comparison uses PBE0/def2-SVP with identical Cartesian
starting geometries.  Gaussian 16 Rev. C.02 native and ESALC-I/\texttt{ReadAllGIC} routes
complete dioxane (4/4), morpholine (6/5), piperazine (6/5), proline (18/16),
ibuprofen (9/9), regenerated uridine (15/15), regenerated tryptophan (9/10),
and artemisinin (10/11) steps, where each pair is listed as native/ESALC-I.
The tryptophan and artemisinin pairs complete the reported
{validation set}.


The same export was compared with Gaussian redundant internals, Cartesians, and
a sequential Z matrix, using \texttt{CalcFC} for ESALC-I.  Water and hydrogen
peroxide are simple controls; norbornane and 75-dimensional camphor test bridged
rings.  The results appear in
Table~\ref{tab:gaussian-coordinate-comparison}. The generated ESALC-I coordinates
reach normal termination in the ordinary and bridged-ring systems.

\begin{table}[htbp]
\centering
\caption{{Gaussian 16 Rev. C.02 HF/STO-3G steps for native redundant,
ESALC-I, Cartesian, and Z-matrix coordinate paths.}}
\label{tab:gaussian-coordinate-comparison}
\begin{tabular}{lccrrrr}
\toprule
System & PG & \ESALCI\ rank & \ESALCI & Default & Cartesian & Z matrix \\
\midrule
Water & \(C_{2v}\) & 3/3 & 5 & 4 & 4 & 4 \\
Hydrogen peroxide & \(D_{2h}\) & 6/6 & 6 & 6 & 6 & 5 \\
Norbornane & \(C_{2v}\) & 51/51 & 2 & 4 & 4 & 6 \\
Camphor & \(C_1\) & 75/75 & 4 & 8 & 14 & 33 \\
\bottomrule
\end{tabular}
\end{table}

A separate compatibility check retains the curvilinear \(U\) representation
throughout.  Camphor contains one \(U(6,4,7,8)\) row and converges in eight
HF/STO-3G steps with the internal evaluator.  Its Gaussian 16 Rev. C.02 input translates
that local source into the supported ordered four-point syntax, while retaining
the original coordinate value, derivative, and family assignment.  SI
Sec.~S3 gives the exact atom-order convention and the frozen rows of the
51-coordinate norbornane example.


\section{Baker transition-state geometries}
\label{sec:baker-geometries}
The following Cartesian starting geometries were used for both the ESALC-I and native Gaussian paths in the 25-case Baker comparison.\cite{BakerChan1996TS} Coordinates are in angstrom; charge and spin multiplicity are given for each case.
\subsection*{Baker TS 01: \texttt{hcn}}
Charge/multiplicity: \texttt{0 1}.
\begin{verbatim}
C 0.000000 0.000000 0.000000
N 0.000000 0.000000 1.148380
H 1.585360 0.000000 1.148380
\end{verbatim}
\subsection*{Baker TS 02: \texttt{hcch}}
Charge/multiplicity: \texttt{0 1}.
\begin{verbatim}
C 0.000000 0.000000 0.000000
C 0.000000 0.000000 1.240540
H -1.440083 0.000000 0.819518
H 0.525846 0.000000 -0.924033
\end{verbatim}
\subsection*{Baker TS 03: \texttt{h2co}}
Charge/multiplicity: \texttt{0 1}.
\begin{verbatim}
C 0.000000 0.000000 0.000000
O 0.000000 0.000000 1.250000
H 1.280250 0.000000 -0.225743
H 0.750000 0.000000 -1.299038
\end{verbatim}
\subsection*{Baker TS 04: \texttt{ch3o}}
Charge/multiplicity: \texttt{0 2}.
\begin{verbatim}
O 0.000000 0.000000 0.000000
C 0.000000 0.000000 1.423000
H 1.006389 0.000000 0.332387
H -0.272228 0.924952 1.924921
H -0.272228 -0.924952 1.924921
\end{verbatim}
\subsection*{Baker TS 05: \texttt{cyclopropyl}}
Charge/multiplicity: \texttt{0 2}.
\begin{verbatim}
C 0.000000 0.000000 0.000000
C 0.000000 0.000000 1.454000
C 1.431910 0.000000 1.201516
H 0.736675 -0.618144 -0.546300
H -0.903666 0.328908 -0.546300
H 2.021514 0.914594 1.399307
H 2.021514 -0.914594 1.399307
H -0.798356 0.096538 2.150710
\end{verbatim}
\subsection*{Baker TS 06: \texttt{bicyclobutane}}
Charge/multiplicity: \texttt{0 1}.
\begin{verbatim}
C 0.000000 0.000000 0.000000
C 0.000000 0.000000 1.495000
C 1.417048 0.000000 1.546961
C -0.654276 1.115186 0.684581
H 0.757024 -0.058250 -0.786238
H -0.605360 -0.701316 2.108202
H 1.994775 -0.884164 1.847108
H 2.044468 0.897275 1.478764
H -0.296120 2.164605 0.735031
H -1.759901 1.143857 0.687196
\end{verbatim}
\subsection*{Baker TS 07: \texttt{bicyclobutane}}
Charge/multiplicity: \texttt{0 1}.
\begin{verbatim}
C 0.000000 0.000000 0.000000
C 0.000000 0.000000 1.489000
C 0.949272 0.000000 2.482357
C -0.106958 1.329369 0.369858
H -0.845880 -0.478576 -0.532088
H -0.981949 -0.388781 1.858846
H 1.963902 0.407430 2.381697
H 0.705633 -0.224075 3.528225
H 0.736421 1.988287 0.648738
H -0.968647 2.008493 0.176780
\end{verbatim}
\subsection*{Baker TS 08: \texttt{formyloxyethyl}}
Charge/multiplicity: \texttt{0 2}.
\begin{verbatim}
C 0.000000 0.000000 0.000000
C 0.000000 0.000000 1.487000
O 1.385059 0.000000 1.942384
C 2.250680 0.190043 0.904208
O 1.694585 0.623486 -0.151624
H -0.523112 0.805522 -0.525845
H 0.067441 -0.964456 -0.516229
H -0.471937 0.916220 1.922864
H -0.439409 -0.926357 1.937758
H 3.266797 -0.190300 1.044500
\end{verbatim}
\subsection*{Baker TS 09: \texttt{parentdieslalder}}
Charge/multiplicity: \texttt{0 1}.
\begin{verbatim}
C 0.0000000000 1.4000000000 0.0000000000
C 0.0000000000 -1.4000000000 0.0000000000
C 1.2000000000 -0.7700000000 0.0000000000
C 1.2000000000 0.7700000000 0.0000000000
C -1.0000000000 0.7000000000 -1.7320510000
C -1.0000000000 -0.7000000000 -1.7320510000
H 2.1112230000 -1.3439580000 0.0815170000
H 2.1112230000 1.3439580000 0.0815170000
H -0.0073520000 2.4767720000 0.0831100000
H -0.0073520000 -2.4767720000 0.0831100000
H -0.8958700000 0.8032440000 -0.0993020000
H -0.8958700000 -0.8032440000 -0.0993020000
H -0.2195930000 1.2065600000 -2.2784970000
H -0.2195930000 -1.2065600000 -2.2784970000
H -1.8670670000 1.2153310000 -1.3460080000
H -1.8670670000 -1.2153310000 -1.3460080000
\end{verbatim}
\subsection*{Baker TS 10: \texttt{tetrazine}}
Charge/multiplicity: \texttt{0 1}.
\begin{verbatim}
N 1.310851 0.631955 0.161807
N 0.425507 1.348536 -0.128337
C 1.194308 -0.762881 0.062226
C -2.016656 0.303884 0.061906
N -0.051060 -1.145422 -0.085710
N -0.912607 -0.213309 -0.200128
H 2.062050 -1.413074 0.105162
H -2.536802 -0.185255 0.916616
\end{verbatim}
\subsection*{Baker TS 11: \texttt{trans\_butadiene}}
Charge/multiplicity: \texttt{0 1}.
\begin{verbatim}
C 0.000000000 0.000000000 0.000000000
C 0.000000000 0.000000000 1.467000000
C 1.094330000 0.000000000 -0.738135000
C -0.703421002 0.838305086 2.205135000
H -0.925741000 0.000000000 -0.556241000
H 0.595054891 -0.709158695 2.023241000
H 1.000114000 0.000000000 -1.815021000
H 2.074955000 0.000000000 -0.280862000
H -0.642861174 0.766131784 3.282021000
H -1.333755653 1.589507444 1.747862000
\end{verbatim}
\subsection*{Baker TS 12: \texttt{ethane\_h2\_abstraction}}
Charge/multiplicity: \texttt{0 1}.
\begin{verbatim}
C 0.000000 0.000000 0.000000
C 0.000000 0.000000 1.430000
H 1.299038 0.000000 -0.750000
H 1.299038 0.000000 0.680000
H -0.238989 -0.958531 -0.460654
H -0.238989 0.958531 -0.460654
H -0.238989 0.958531 1.890654
H -0.238989 -0.958531 1.890654
\end{verbatim}
\subsection*{Baker TS 13: \texttt{hf\_abstraction}}
Charge/multiplicity: \texttt{0 1}.
\begin{verbatim}
C 0.0000000000 0.0000000000 0.0000000000
C 0.0000000000 0.0000000000 1.4300000000
H 1.4772120000 0.0000000000 -0.2604720000
F 1.8711350000 0.0000000000 1.7599320000
H -0.2533350000 -0.9454600000 -0.4564280000
H -0.2533350000 0.9454600000 -0.4564280000
H -0.2533350000 0.9454600000 1.8864280000
H -0.2533350000 -0.9454600000 1.8864280000
\end{verbatim}
\subsection*{Baker TS 14: \texttt{vinyl\_alcohol}}
Charge/multiplicity: \texttt{0 1}.
\begin{verbatim}
C 0.0000000000 0.0000000000 0.0000000000
C 0.0000000000 0.0000000000 1.4300000000
O 1.1258330000 0.0000000000 2.0800000000
H -0.4328730000 -0.9282990000 -0.3728020000
H -0.5121320000 0.8870390000 -0.3728020000
H -0.9403760000 0.0822720000 1.9750000000
H 1.1236090000 0.1018970000 0.7338530000
\end{verbatim}
\subsection*{Baker TS 15: \texttt{hocl}}
Charge/multiplicity: \texttt{0 1}.
\begin{verbatim}
O 0.0000000000 0.0000000000 0.0000000000
C 0.0000000000 0.0000000000 1.1700000000
Cl 2.3350000000 0.0000000000 1.1700000000
H 0.0000000000 0.0000000000 2.2970000000
\end{verbatim}
\subsection*{Baker TS 16: \texttt{h2po4\_anion}}
Charge/multiplicity: \texttt{-1 1}.
\begin{verbatim}
P 0.0000000 0.0000000 0.0000000
O 0.0000000 0.0000000 1.4600000
O 1.3765120 0.0000000 -0.4866370
O -0.9333880 -1.6166750 -0.6599600
O -1.2499610 0.1678910 -0.8837480
H -1.2421250 -2.1282930 0.0785710
H -1.8181770 -0.9580920 -1.0327190
\end{verbatim}
\subsection*{Baker TS 17: \texttt{claisen}}
Charge/multiplicity: \texttt{0 1}.
\begin{verbatim}
C 0.000000 0.000000 0.000000
C 0.000000 0.000000 1.420000
C 1.216766 0.000000 2.122500
O 1.939676 -1.811097 1.643551
C 1.969182 -1.792073 0.334021
C 0.764394 -1.764457 -0.397977
H 0.855525 0.493938 -0.460654
H -0.987875 0.000000 -0.460654
H -0.943968 0.000000 1.965000
H 1.121766 0.000000 3.208352
H 2.043466 0.493938 1.611920
H 0.883818 -1.750040 -1.481319
H -0.080384 -2.253943 0.086640
H 2.925095 -1.797591 -0.189716
\end{verbatim}
\subsection*{Baker TS 18: \texttt{silyene\_insertion}}
Charge/multiplicity: \texttt{0 1}.
\begin{verbatim}
C 0.000000 0.000000 0.000000
C 0.000000 0.000000 1.520000
Si 2.006545 0.000000 2.644179
H 0.422061 0.910502 -0.450390
H 0.459439 -0.873991 -0.461688
H -1.051651 -0.019825 -0.354596
H -0.547586 0.879391 1.858999
H -0.559830 -0.877067 1.876854
H 2.812789 -1.154016 2.187433
H 1.220827 0.893189 1.730472
H 1.749593 -0.143498 4.094622
\end{verbatim}
\subsection*{Baker TS 19: \texttt{hnccs}}
Charge/multiplicity: \texttt{0 1}.
\begin{verbatim}
H 0.0000000000 0.0000000000 0.0000000000
N 0.0000000000 0.0000000000 1.0100000000
C 0.6710310000 0.0000000000 2.0048450000
C 2.5252980000 0.0000000000 2.1507790000
S 3.3197710000 0.0000000000 3.5049250000
\end{verbatim}
\subsection*{Baker TS 20: \texttt{hconh3\_cation}}
Charge/multiplicity: \texttt{1 1}.
\begin{verbatim}
N 0.0000000 0.0000000 0.0000000
C 0.0000000 0.0000000 2.0000000
O 0.9786090 0.0000000 2.5650000
H 0.9709210 0.0000000 -0.3438210
H -0.9353070 0.0000000 1.4600000
H -0.4854600 -0.8408420 -0.3438210
H -0.4854600 0.8408420 -0.3438210
\end{verbatim}
\subsection*{Baker TS 21: \texttt{acrolein\_rot}}
Charge/multiplicity: \texttt{0 1}.
\begin{verbatim}
C 0.000000 0.000000 0.000000
C 0.000000 0.000000 1.340000
C 1.255737 0.000000 2.065000
O 1.784012 1.056551 2.370000
H -0.935307 0.000000 -0.540000
H 0.935307 0.000000 -0.540000
H -0.935307 0.000000 1.880000
H 1.723391 -0.935307 2.335000
\end{verbatim}
\subsection*{Baker TS 22: \texttt{hconhoh}}
Charge/multiplicity: \texttt{0 1}.
\begin{verbatim}
O 0.000000 0.000000 0.000000
C 0.000000 0.000000 1.300000
N 1.221600 0.000000 1.744626
O 1.464708 0.000000 3.123357
H -0.842649 0.000000 2.007066
H 2.450704 0.000000 3.290126
H 1.352334 0.000000 0.250334
\end{verbatim}
\subsection*{Baker TS 23: \texttt{hcn\_h2}}
Charge/multiplicity: \texttt{0 1}.
\begin{verbatim}
H 0.0000000000 0.0000000000 0.0000000000
N 0.0000000000 0.0000000000 1.0000000000
C 1.0392300000 0.0000000000 1.6000000000
H 1.5430290000 0.0868240000 2.4594470000
H 0.5092680000 0.0593310000 3.0682030000
\end{verbatim}
\subsection*{Baker TS 24: \texttt{h2cnh}}
Charge/multiplicity: \texttt{0 1}.
\begin{verbatim}
H 0.0000000000 0.0000000000 0.0000000000
N 0.0000000000 0.0000000000 1.0000000000
C 1.0392300000 0.0000000000 1.6000000000
H 0.1948560000 -0.9750000000 1.7625000000
H 1.9052560000 0.0000000000 1.1000000000
\end{verbatim}
\subsection*{Baker TS 25: \texttt{hcnh2}}
Charge/multiplicity: \texttt{0 1}.
\begin{verbatim}
C 0.000000 0.000000 0.000000
N 0.000000 0.000000 1.350000
H 0.965926 0.000000 -0.258819
H -1.171123 0.676148 1.712347
H -1.236231 -0.449951 1.828828
\end{verbatim}
}

\bibliographystyle{achemso}
\bibliography{archive/original_submission/references}

\end{document}
